\documentclass[11pt]{article}

\usepackage[a4paper,margin=1in]{geometry}
\usepackage{amsmath,amssymb,mathtools}
\usepackage{booktabs}
\usepackage{array}
\usepackage{enumitem}
\usepackage{algorithm}
\usepackage{algpseudocode}
\usepackage{xcolor}
\usepackage{hyperref}
\usepackage{float}
\usepackage{graphicx}

\title{fActorI: A Multi-Fidelity Variance-Driven Agentic Framework for Scientific Discovery}
\author{Architecture Design Document}
\date{}

\newcommand{\R}{\mathbb{R}}
\newcommand{\Cands}{\mathcal{C}}
\newcommand{\Tree}{\mathcal{T}}
\newcommand{\Papers}{\mathcal{P}}
\newcommand{\Models}{\mathcal{M}}
\newcommand{\Reviewers}{\mathcal{R}}
\newcommand{\Fields}{\mathcal{F}}
\newcommand{\SysState}{\mathcal{S}}
\newcommand{\clip}{\operatorname{clip}}
\newcommand{\cosim}{\operatorname{cos}}
\newcommand{\argtop}{\operatorname{argtop}}
\newcommand{\mean}{\operatorname{mean}}
\newcommand{\Var}{\operatorname{Var}}
\newcommand{\Cost}{\operatorname{Cost}}
\newcommand{\freq}{\operatorname{freq}}
\newcommand{\Hash}{\operatorname{Hash}}
\newcommand{\Diff}{\operatorname{Diff}}
\newcommand{\Merge}{\operatorname{Merge}}

\begin{document}

\maketitle

\section{Goal}

The primary goal of the system is to support and partially automate the scientific research pipeline from opportunity discovery to formally labeled manuscript synthesis. It conceptualizes research discovery as a \emph{lift-then-filter} process. First, the system lifts low-dimensional user constraints into a high-dimensional research space by generating diverse, mathematically structured candidate branches. Then, through staged validation, adversarial review, and verification-aware pruning, it filters this expanded space into a smaller set of viable research directions.

The tool operates as a multi-fidelity, variance-driven scientific paper writer agent. It deliberately avoids converging too early on a single answer. Instead, it explores a tree of candidate research programs, maintains diversity through ensemble distance metrics, and removes semantic redundancy through approximate branch-distance constraints. The purpose of this variance is not to generate more text, but to expose alternative hypotheses, methods, mathematical lenses, baselines, and verification paths before expensive validation begins.

The intended operational mode is \emph{factory-like}: the system should autonomously generate, repair, validate, downgrade, reject, or synthesize research branches, while reserving human intervention for tail cases involving high uncertainty, high risk, extra budget, external access, irreversible decisions, or authorial preference. Ordinary decisions such as narrowing a question, adding data, strengthening a baseline, downgrading a claim, stopping a failed branch, or selecting a clearly dominant final nucleus are handled by bounded controllers rather than by repeated human approval.

To minimize compute footprint, the system applies a staged cheap-to-expensive validation funnel. Early stages use low-cost opportunity scoring, question refinement, literature checks, duplicate filtering, and bridge-risk analysis. Later stages allocate more expensive resources only to promising branches, including adversarial red-team review, Lean-based mathematical verification, empirical experiment controllers, and final synthesis.

The system enforces strict epistemic separation between exploration, verification, and presentation. Candidate ideas may be explored speculatively, but final theorem claims require Lean verification, final empirical claims require predeclared experimental success, and all unresolved claims must remain explicitly labeled as conjectures, limitations, negative results, or unsupported directions. LaTeX is treated as a presentation format, not as evidence.

Crucially, the system includes a mandatory opportunity-discovery layer and adaptive state controllers for literature, question refinement, proof, experiment design, bridge validation, disagreement resolution, and budget allocation. These controllers diagnose, refine, repair, pause, or prune branches throughout the execution pipeline. Every state-changing action is recorded in an immutable hash-linked research ledger, with artifacts stored separately and linked by content hashes. The ledger provides a reproducible audit trail of decisions, prompts, scores, verification attempts, failures, and final claims.

The final convergence stage does not merely select the highest-scoring surviving branch. When several verified or partially verified branches reveal a common structure, the system attempts to synthesize a higher-level abstraction---a general model, framework, theorem family, or empirical mechanism---under which the surviving branches become special cases. If no valid abstraction survives adversarial scrutiny, the system falls back to a single best verified branch.

The intended output is therefore not merely a polished manuscript. It is a ledgered research object: a verified or explicitly labeled paper draft together with its candidate tree, failed branches, proof and experiment artifacts, literature evidence, red-team objections, and provenance trail. For the initial factory-mode implementation, the system restricts empirical branches to paper ideas requiring no data or at most synthetic data, deferring branches that require public or user-provided real datasets.


\begin{figure}[H]
    \centering
    \includegraphics[width=1\linewidth]{Analogy.png}
    \caption{Analogy: Research discovery as a lift-then-project operation. The system maps low-dimensional prompts into an expanded latent idea space (the lift), which is then orthogonally reduced to promising, verified research directions via constraint-based gating (the projection).}
    \label{fig:analogy}
\end{figure}

\section{System Architecture Diagram}

Figure \ref{fig:architecture} illustrates the multi-fidelity execution flow. The architecture fixes the structure of the decision process, detailing the transitions from Stage 0 Opportunity Discovery through the Adaptive Controllers and Red Team Validation Layer, culminating in Stage C Deep Verification.

\begin{figure}[H]
    \centering
    \includegraphics[width=1\linewidth]{ChatGPT Image Jun 14, 2026, 12_58_57 PM.png}
    \caption{Architecture summary}
    \label{fig:architecture}
\end{figure}

\begin{figure}[H]
    \centering
    \includegraphics[width=1\linewidth]{ChatGPT Image Jun 14, 2026, 12_59_33 PM.png}
    \caption{Example research tree}
    \label{fig:research_tree}
\end{figure}

The system executes a generalized constraint-based tree search over the research program space. A user's prompt is parsed into a constraint set $C$, which may fix any subset of the research variables:
$$
D,\quad P,\quad M,\quad Q,\quad H,\quad T,\quad E.
$$
Here $D$ is the domain, $P$ the extracted domain primitives, $M$ the mathematical or computational method, $Q$ the research question, $H$ the hypothesis, $T$ the theoretical structure, and $E$ the experimental contract.

A candidate research program $c$ is admissible only if it satisfies the user's constraints:
$$
c \models C.
$$
Unspecified fields are treated as open search variables. If a user provides only a broad domain constraint $D$, for example $D=\text{``human geography''}$, the system must execute Stage 0 Opportunity Discovery before generating research questions.

\subsection{Preset Interpretations}

The same search engine supports several user input types:
\begin{center}
\begin{tabular}{lll}
\toprule
Input Type & Fixed Variables & Open Variables \\
\midrule
Exact question & $Q$ & $M,H,T,E$ \\
Domain plus method & $D,M$ & $Q,H,T,E$ \\
Domain only & $D$ & $P,M,Q,H,T,E$ \\
Experiment request & $D,M,Q,E$ partially & $H,B,V$ \\
Mathematical theorem request & $D,M,Q,T$ partially & $H,V$ \\
\bottomrule
\end{tabular}
\end{center}

\section{Candidate Representation and Symbolic Mapping}

Each candidate branch is represented as
$$
c = (D, P, M, Q, H, T, E, B, \delta_{\text{data}}, L, V, R, A, \sigma_{\text{sym}}),
$$
where:
\begin{center}
\begin{tabular}{lll}
\toprule
Symbol & Field & Structural Meaning \\
\midrule
$D$ & domain & Broad scientific or academic field \\
$P$ & primitives & Extracted domain structures such as flows, clustering, distance, dependence \\
$M$ & method & Mathematical or computational analytical machinery \\
$Q$ & question & Explicit research question statement \\
$H$ & hypothesis & Central falsifiable claim to investigate \\
$T$ & theory & Mathematical framework: definitions, lemmas, assumptions \\
$E$ & experiment & Pre-committed experimental evaluation contract \\
$B$ & baseline & Reference theorem, model, or state-of-the-art method \\
$\delta_{\text{data}}$ & data regime & Required data access level: no data, synthetic data, public download, or user-provided data \\
$L$ & literature & Retrieved contextually relevant prior works \\
$V$ & verification & Multi-fidelity proof or empirical execution logs \\
$R$ & reviews & Multi-agent feedback and structured critique scores \\
$A$ & artifact & Progressively generated manuscript, proof, or code artifacts \\
$\sigma_{\text{sym}}$ & symbolic state & Extracted meta-features for ensemble distance matching \\
\bottomrule
\end{tabular}
\end{center}

The symbolic state $\sigma_{\text{sym}}$ includes:
\begin{itemize}[leftmargin=*]
\item \textbf{Objects:} mathematical spaces, manifolds, graphs, distributions, metrics, bundles, networks, flows.
\item \textbf{Assumptions:} compactness, smoothness, completeness, sparsity, stationarity, independence, boundedness.
\item \textbf{Claim type:} existence, uniqueness, stability, classification, asymptotic estimate, counterexample, algorithmic guarantee.
\item \textbf{Method class:} optimal transport, copulas, graph theory, PDEs, Ricci flow, causal inference, stochastic processes, topological data analysis.
\item \textbf{Target directive:} theorem, conjecture, empirical model, benchmark, dataset, negative result, survey-to-theorem conversion.
\end{itemize}

\section{Ensemble Distance Metrics}

Candidate deduplication and diversity estimation use an ensemble distance metric:
$$
D_{\text{total}}(c_i,c_j) = 0.35D_{\text{embedding}} + 0.25D_{\text{symbolic}} + 0.20D_{\text{citation}} + 0.20D_{\text{reviewer}}.
$$

For formal mathematical tracks:
$$
D_{\text{total\_math}}(c_i,c_j) = 0.25D_{\text{embedding}} + 0.40D_{\text{symbolic}} + 0.20D_{\text{citation}} + 0.15D_{\text{reviewer}}.
$$

The active branch distance is
$$
D_{\text{track}}(c_i,c_j)=
\begin{cases}
D_{\text{total\_math}}(c_i,c_j), & \text{if } c_i,c_j \text{ are both mathematical},\\
D_{\text{total}}(c_i,c_j), & \text{otherwise}.
\end{cases}
$$

Two branches are duplicates if
$$
D_{\text{track}}(c_i,c_j)\leq \tau_{\text{dup}}, \qquad \tau_{\text{dup}}=0.15.
$$
The lower-scoring branch is pruned.

\subsection{Approximate Orthogonality of Candidate Branches}

The framework does not assume that research ideas are truly orthogonal in a mathematical Hilbert space. Research candidates are structured semantic objects, not native vectors with exact inner products. Instead, the system enforces approximate orthogonality relative to the ensemble candidate metric $D_{\text{track}}$.

Let $K_{\text{track}}$ denote the induced similarity:
$$
K_{\text{track}}(c_i,c_j)=1-D_{\text{track}}(c_i,c_j).
$$

Two candidates are $\epsilon$-orthogonal if:
$$
\operatorname{Orth}_{\epsilon}(c_i,c_j)=1
\Longleftrightarrow
K_{\text{track}}(c_i,c_j)\leq \epsilon.
$$

Equivalently:
$$
\operatorname{Orth}_{\epsilon}(c_i,c_j)=1
\Longleftrightarrow
D_{\text{track}}(c_i,c_j)\geq 1-\epsilon.
$$

A surviving candidate set $\mathcal{S}$ is $\epsilon$-orthogonal if:
$$
\forall c_i,c_j\in\mathcal{S},\ i\neq j,
\qquad
K_{\text{track}}(c_i,c_j)\leq \epsilon.
$$

This condition is not interpreted as absolute geometric orthogonality. It is interpreted as enforced diversity across semantic meaning, symbolic structure, citation neighborhood, and reviewer judgment.

To select diverse branches, the system uses diversity-constrained greedy selection. Given a set of already selected branches $\mathcal{S}$, the next candidate is:
$$
c^\star = \arg\max_{c\in\mathcal{C}\setminus\mathcal{S}}
\left[
S_{\text{cost}}(c)
+
\lambda_{\text{div}}
\min_{c_j\in\mathcal{S}}
D_{\text{track}}(c,c_j)
\right],
$$
where $\lambda_{\text{div}}\geq0$ controls the strength of the diversity pressure.

In addition to duplicate pruning, the system may cluster candidates under $D_{\text{track}}$:
$$
\mathcal{K} = \operatorname{Cluster}(\mathcal{C},D_{\text{track}}).
$$

To avoid over-exploiting one region of idea space, the system enforces a cluster-level survivor cap:
$$
|\mathcal{S}\cap K_j|\leq m_{\text{cluster}}
$$
for each cluster $K_j\in\mathcal{K}$.

Thus, candidate selection is not purely score-maximizing. It is constrained by approximate orthogonality:
$$
\text{high score}
+
\text{low redundancy}
+
\text{coverage of distinct research directions}.
$$

This prevents the search tree from collapsing into many near-identical variants of the same idea.

\section{Base Component Scoring}

Each candidate receives structural scores in $[0,1]$.

\subsection{Novelty Score}

$$
N(c) = 0.55N_{\text{lit}}(c) + 0.25N_{\text{set}}(c) + 0.20N_{\text{comb}}(c).
$$

Here:
\begin{itemize}[leftmargin=*]
\item $N_{\text{lit}}(c)$ measures distance from retrieved prior literature.
\item $N_{\text{set}}(c)$ measures distance from already accepted branches.
\item $N_{\text{comb}}(c)$ measures novelty of combining distant literatures or primitives.
\end{itemize}

\subsection{Feasibility Score}

$$
F(c) = 0.20F_{\text{resources}} + 0.20F_{\text{baseline}} + 0.20F_{\text{tooling}} + 0.20F_{\text{time}} + 0.20F_{\text{scope}}.
$$

\subsection{Verifiability Score}

$$
V(c) = 0.30V_{\text{hyp}} + 0.25V_{\text{metric}} + 0.20V_{\text{baseline}} + 0.15V_{\text{failure}} + 0.10V_{\text{repro}}.
$$

\subsection{Difficulty Score}

$$
\operatorname{Diff}(c) = 0.25D_{\text{def}}(c) + 0.25D_{\text{dep}}(c) + 0.25D_{\text{search}}(c) + 0.25D_{\text{cost}}(c).
$$

The default admissible difficulty interval for autonomous discovery is
$$
0.35\leq \operatorname{Diff}(c)\leq 0.70.
$$

\subsection{Base Candidate Score}

$$
S(c) = 0.25N(c) + 0.20F(c) + 0.20V(c) + 0.15R(c) + 0.10(1-\operatorname{Diff}(c)) + 0.10D_{\text{rev}}(c).
$$

\section{Stage 0: Opportunity Discovery}

Stage 0 prevents the system from remaining inside conventional domain vocabulary. It is mandatory whenever the user specifies a broad domain $D$ but does not specify a mathematical method $M$.

Given $D$, the system extracts primitives $P(D)=\{p_1,\dots,p_k\}$. For example, if $D=\text{``human geography''}$, possible primitives include: flows, distance, migration, commuting, accessibility, spatial inequality, regional dependence, urban networks, mobility, segregation.

\subsection{Mathematical Method Candidate Set}

The candidate mathematical method set is $\mathcal{M}_{\text{cand}} = \mathcal{M}_{\text{library}} \cup \mathcal{M}_{\text{retrieved}} \cup \mathcal{M}_{\text{generated}}.$ Examples include: optimal transport, copulas, graph curvature, spatial statistics, stochastic processes, topological data analysis, causal inference, information geometry.

\subsection{Easy-Win Opportunity Score}

Each domain-method pair $(D,M_i)$ receives an opportunity score:
$$
O(D,M_i) = 0.25S_{\text{fit}} + 0.20S_{\text{underuse}} + 0.20S_{\text{abundance}} + 0.15S_{\text{verify}} + 0.10S_{\text{data}} + 0.10S_{\text{paper}}.
$$

\subsection{Quantitative Underuse Score and Evidence Attestation}

Let $\mathcal{L}_{\text{math}}$ be a broad mathematical literature corpus and $\mathcal{L}_{D}$ be the domain-specific literature corpus. The system calculates empirical underuse:
$$
S_{\text{underuse\_raw}}(D,M) = \clip\left( \frac{\freq(M,\mathcal{L}_{\text{math}}) - \freq(M,\mathcal{L}_{D})}{\freq(M,\mathcal{L}_{\text{math}})+\epsilon}, 0, 1 \right), \qquad \epsilon=10^{-8}.
$$

To prevent opportunity hallucination, the system explicitly splits underuse into $\text{Underuse}_{\text{good}}$ and $\text{Underuse}_{\text{bad}}$. Underuse contributes positively to the opportunity score only if structural fit and verification feasibility are high:
$$
S_{\text{underuse}}(D,M) = \begin{cases} S_{\text{underuse\_raw}}(D,M), & \text{if } S_{\text{fit}}(D,M) \geq 0.70 \text{ and } S_{\text{verify}}(D,M) \geq 0.60 \\ 0, & \text{otherwise} \end{cases}
$$

The final opportunity score $O_{\text{final}}(D,M)$ applies a False-Bridge Penalty $S_{\text{forced}}(D,M)$:
$$
O_{\text{final}}(D,M) = O(D,M) - 0.20S_{\text{forced}}(D,M).
$$

The system promotes a mathematical lens only if $O_{\text{final}}(D,M)\geq 0.70$. Promoted methods seed the constraint sets $\mathcal{C}_{0} = \{C\cup(M=M_i): O_{\text{final}}(D,M_i)\geq0.70\}$.

\section{Research Density Landscape}

The system may be interpreted as searching over a latent space of admissible research programs. Each candidate paper is represented by a structured minimal description:
$$ x(c) = (D,P,M,Q,H,T,E,B), $$
where $D$ is the domain, $P$ the extracted primitives, $M$ the method, $Q$ the question, $H$ the hypothesis, $T$ the theoretical structure, $E$ the empirical contract, and $B$ the baseline.

Let $\mathcal{X}_C$ denote the space of admissible paper descriptions satisfying the user constraint set $C$:
$$ \mathcal{X}_C = \{x(c): c\models C\}. $$

The purpose of the search is not merely to find isolated high-scoring ideas, but to identify regions of $\mathcal{X}_C$ that are simultaneously underexploited, feasible, verifiable, and scientifically meaningful.

\begin{figure}[H]
    \centering
    \includegraphics[width=0.85\linewidth]{Research_density_landscape.png}
    \caption{The Latent Research Density Landscape. The search algorithm navigates away from crowded literature regions towards underexploited, feasible basins.}
    \label{fig:density_landscape}
\end{figure}

\subsection{Embedding of Minimal Paper Descriptions}
Each minimal description is mapped into a latent research space:
$$ \phi:\mathcal{X}_C\rightarrow \mathbb{R}^d. $$

The embedding $\phi(x)$ is not a single text embedding. It is an ensemble representation:
$$ \phi(x) = \phi_{\text{ens}}(D,P,M,Q,H,T,E,B), $$
combining semantic embeddings, symbolic structure, citation context, method class, claim type, and verification target.

A distance on admissible paper descriptions is defined by:
$$ d_{\mathcal{X}}(x_i,x_j) = D_{\text{track}}(c_i,c_j), $$
where $D_{\text{track}}$ is the ensemble branch distance used for duplicate filtering and diversity control.

\subsection{Prior-Work Density}
Let $\mathcal{P}$ be the retrieved literature corpus relevant to the user constraints. Each prior paper $p\in\mathcal{P}$ is mapped to the same latent space:
$$ z_p=\phi_{\text{lit}}(p)\in\mathbb{R}^d. $$

The local prior-work density around a candidate $x$ is estimated by kernel density:
$$ \rho_{\text{lit}}(x) = \sum_{p\in\mathcal{P}} K_{\sigma} \left( d_{\mathcal{X}}(x,p) \right), $$
where:
$$ K_{\sigma}(r) = \exp\left(-\frac{r^2}{2\sigma^2}\right). $$

High $\rho_{\text{lit}}(x)$ means the candidate lies in a crowded region of existing literature. Low $\rho_{\text{lit}}(x)$ means the candidate lies in a sparse or underexplored region.

To prevent artificial novelty caused by retrieval failure, the density estimate is admissible only if the literature controller satisfies:
$$ \rho_{\text{adequacy}}(x)\geq \tau_{\text{adequacy}}. $$
If this condition fails, the density estimate is labeled unreliable:
$$ \operatorname{DensityReliable}(x)=0. $$

\subsection{Difficulty-Weighted Research Potential}
Sparse regions are not automatically valuable. Some regions are sparse because they are impossible, irrelevant, unverifiable, or computationally infeasible. Therefore the system defines a difficulty-weighted density objective.

Let $\operatorname{Diff}(x)\in[0,1]$ be the candidate difficulty score, and let $I(x), F(x), V(x)$ denote importance, feasibility, and verifiability respectively.

The underexploitation energy is defined as:
$$ \mathcal{E}(x) = \rho_{\text{lit}}(x) + \lambda_{\text{diff}}\operatorname{Diff}(x) + \lambda_{\text{risk}}R_{\text{risk}}(x) - \lambda_{\text{val}} \left( I(x)F(x)V(x) \right), $$
where $R_{\text{risk}}(x)$ aggregates bridge risk, novelty risk, baseline risk, and verification risk.

The system searches for local minima of $\mathcal{E}$:
$$ x^\star \in \operatorname*{arg\,locmin}_{x\in\mathcal{X}_C} \mathcal{E}(x). $$

A low value of $\mathcal{E}(x)$ corresponds to a candidate that lies in a sparse region of prior work, has manageable difficulty, avoids known risk modes, and retains scientific value.

\subsection{Underexploited Feasible Basin}
A region $U\subset\mathcal{X}_C$ is called an underexploited feasible basin if:
$$ \mathbb{E}_{x\in U}[\rho_{\text{lit}}(x)]\leq \tau_{\rho}, $$
$$ \mathbb{E}_{x\in U}[\operatorname{Diff}(x)]\in[\tau_{\text{diff,min}},\tau_{\text{diff,max}}], $$
and
$$ \mathbb{E}_{x\in U}[F(x)V(x)]\geq \tau_{\text{fv}}. $$

Thus, the system does not merely search for isolated novelty. It searches for neighborhoods of research space that appear systematically underdeveloped while remaining feasible and verifiable.

\subsection{Relation to Opportunity Discovery}
The Stage 0 opportunity score can be interpreted as an approximate basin-finding heuristic:
$$ O_{\text{final}}(D,M) \approx \operatorname{BasinScore}(U_{D,M}), $$
where $U_{D,M}$ is the subset of admissible paper descriptions using domain $D$ and method $M$:
$$ U_{D,M} = \{x\in\mathcal{X}_C:x_D=D,\ x_M=M\}. $$

The basin score is:
$$ \operatorname{BasinScore}(U) = \alpha_1(1-\bar{\rho}_{\text{lit}}(U)) + \alpha_2\bar{F}(U) + \alpha_3\bar{V}(U) + \alpha_4\bar{I}(U) - \alpha_5\overline{\operatorname{Diff}}(U) - \alpha_6\bar{R}_{\text{risk}}(U). $$
This connects the opportunity-discovery layer to a geometric interpretation: Stage 0 searches for low-density, high-potential regions of the admissible research landscape.

\subsection{Local Search Dynamics}
The research tree can be viewed as a local search process over $\mathcal{X}_C$. Each controller action produces a neighboring candidate:
$$ x_{t+1} = a(x_t), $$
where
$$ 
\begin{aligned}
a \in \{ &\textsc{NarrowScope}, \textsc{ChangeMetric}, \textsc{AddBaseline}, \textsc{SearchAdjacentFields}, \\
&\textsc{DefineObjects}, \textsc{SplitLemma}, \textsc{AddDataset}, \textsc{StrengthenBaseline} \}.
\end{aligned}
$$

The controller accepts a refinement if it improves the conservative objective:
$$ \mathcal{E}(x_{t+1})+\eta_t < \mathcal{E}(x_t), $$
where $\eta_t$ is a tolerance term allowing controlled exploration when reviewer disagreement, novelty, or uncertainty suggests that a nearby branch may be valuable.

The search therefore behaves as a constrained descent process over the research density landscape, with occasional diversity-preserving moves to avoid premature convergence.

\subsection{Pruning Rule}
A candidate is pruned if it lies in a crowded, high-risk, or overly difficult region:
$$ \operatorname{Prune}(x)=1 $$
if any of the following hold:
$$ \rho_{\text{lit}}(x)>\tau_{\rho,\max}, $$
$$ \operatorname{Diff}(x)>\tau_{\text{diff,max}}, $$
$$ F(x)V(x)<\tau_{\text{fv}}, $$
$$ R_{\text{risk}}(x)>\tau_{\text{risk}}, $$
or
$$ \hat{S}(x)-u_S(x)<\tau_S. $$

A candidate may survive despite high difficulty only if it is explicitly preserved by the budget scheduler or human operator and remains labeled as high-risk until verification succeeds.

\subsection{Interpretation}
This formulation makes the search objective explicit. The system is not attempting to maximize paper production volume. It is attempting to locate underexploited feasible basins in the admissible research landscape:
$$ \text{good research direction} \approx \text{low prior-work density} + \text{manageable difficulty} + \text{high verifiability} + \text{scientific value}. $$

The final manuscript is then the projection of a surviving basin candidate into a labeled scientific artifact, with all verification status, failed branches, and provenance retained in the ledger.

\section{Adaptive Question Refiner}

Following Stage 0 generation, questions are often structurally imprecise. The system executes an adaptive refinement loop to optimize the problem framing:
Let the question state be:
$$
Q_{\text{state}}(c) = (S_{\text{scope}}, S_{\text{novelty}}, S_{\text{testable}}).
$$

The controller selects an action:
$$
\begin{aligned}
a_Q(c) \in \{ & \textsc{NarrowScope}, \textsc{BroadenScope}, \textsc{ChangeObject}, \textsc{ChangeMetric}, \\
& \textsc{AddBaseline}, \textsc{ConvertToConjecture}, \textsc{ConvertToExperiment}, \textsc{RejectQuestion} \}
\end{aligned}
$$

\section{Adaptive Literature Controller and Retrieval Adequacy Certificate}

Retrieval cannot be fixed at $K$ papers without risking undetected redundancy, missed terminology, or false novelty. The system therefore replaces the notion of absolute literature coverage with a \emph{retrieval adequacy certificate}. The certificate does not claim that the entire literature has been exhausted. It certifies that several complementary retrieval modes have been executed and that the remaining novelty risk is low enough for the current stage.

Let the literature state be:
$$
L_{\text{state}}(c) = (K, \rho_{\text{semantic}}, \rho_{\text{keyword}}, \rho_{\text{citation}}, \rho_{\text{diversity}}, \rho_{\text{adversarial}}, \rho_{\text{novelty-risk}}).
$$

The retrieval adequacy score is:
$$
\rho_{\text{adequacy}}(c)
=
w_1\rho_{\text{semantic}}
+w_2\rho_{\text{keyword}}
+w_3\rho_{\text{citation}}
+w_4\rho_{\text{diversity}}
+w_5\rho_{\text{adversarial}},
$$
with $w_i\geq0$ and $\sum_i w_i=1$.

The components have the following interpretation:
\begin{center}
\begin{tabular}{p{0.24\linewidth} p{0.64\linewidth}}
\toprule
Component & Meaning \\
\midrule
$\rho_{\text{semantic}}$ & Dense retrieval over candidate descriptions, abstracts, and related concepts. \\
$\rho_{\text{keyword}}$ & Boolean and lexical retrieval using explicit terminology, synonyms, and negative controls. \\
$\rho_{\text{citation}}$ & Citation-neighborhood traversal from seed papers, surveys, and closest priors. \\
$\rho_{\text{diversity}}$ & Diversity across subdomains, time periods, methodological schools, and terminology clusters. \\
$\rho_{\text{adversarial}}$ & Targeted search for prior work likely to invalidate novelty or reveal equivalent formulations. \\
\bottomrule
\end{tabular}
\end{center}

The adaptive controller selects an action based on diagnostic metrics:
$$
\begin{aligned}
a_{\text{lit}}(c) \in \{ & \textsc{StopRetrieval}, \textsc{ExpandKeywords}, \textsc{SearchAdjacentFields}, \\
& \textsc{SearchAdversarialNovelty}, \textsc{IncreaseK}, \textsc{TraverseCitationGraph}, \\
& \textsc{QuerySurveyPapers}, \textsc{RejectKnown} \}.
\end{aligned}
$$

\textbf{Stop Condition.} The controller may return $\textsc{StopRetrieval}$ only if:
$$
\rho_{\text{adequacy}}(c)\geq \tau_{\text{adequacy}}
\quad \text{and} \quad
\rho_{\text{novelty-risk}}(c)\leq \tau_{\text{novelty-risk}}.
$$
The default MVP thresholds are:
$$
\tau_{\text{adequacy}}=0.80,
\qquad
\tau_{\text{novelty-risk}}=0.20.
$$

To prevent artificial novelty via retrieval failure, Stage C execution is strictly forbidden when:
$$
\rho_{\text{adequacy}}(c)<\tau_{\text{adequacy}}.
$$
Such a branch is labeled:
$$
\textsc{InsufficientRetrievalAdequacy}.
$$

The certificate is stored as an artifact and committed to the ledger. It must include the exact queries, retrieved papers, citation traversals, closest-prior comparisons, rejected known-prior matches, and unresolved blindspots.

\section{Three-Stage Multi-Fidelity Funnel}

The framework scales search granularity over a multi-fidelity funnel, aggressively pruning search spaces prior to committing resource-heavy evaluation.

\begin{table}[H]
\centering
\caption{Funnel Configuration Matrix}
\begin{tabular}{p{0.22\linewidth} p{0.22\linewidth} p{0.22\linewidth} p{0.22\linewidth}}
\toprule
Parameter & Stage A: Cheap Search & Stage B: Mid Validation & Stage C: Full Verification \\
\midrule
Branches Generated & 12 & 4 per active branch & -- \\
Surviving Targets & Top 4 & Top 2 & Final candidate \\
Literature Search & Adaptive ($K\approx5$) & Adaptive ($K\approx20$) & Adaptive ($K\approx50$) \\
Reviewer Pool & 1 critic & 3 expert reviewers & 5 adversarial reviewers \\
Lean Gating Level & Level 0--1 & Level 2--3 & Level 4--5 \\
Empirical Testing & Static contract check & Toy or small-scale proxy & Full contract deployment \\
\bottomrule
\end{tabular}
\end{table}

\section{Adaptive Budget Scheduler}

Fixed scalar budgets are safe but sub-optimal. The Adaptive Budget Scheduler treats compute boundaries as a strict multi-dimensional resource vector:
$$
\mathbf{B} = (B_{\text{api}}, B_{\text{retrieval}}, B_{\text{lean}}, B_{\text{gpu}}).
$$

The system enforces independent non-negative hard caps across all modalities:
$$
\mathbf{B}_{\text{spent}} \leq \mathbf{B}_{\max}, \quad \text{i.e.,} \quad C_{\text{llm}}\leq B_{\text{api\_max}}, \ C_{\text{retrieval}}\leq B_{\text{retrieval\_max}}, \ C_{\text{lean}}\leq B_{\text{lean\_max}}, \ C_{\text{experiment}}\leq B_{\text{gpu\_max}}.
$$

For ranking, costs are normalized into a unified scalar $\widehat{\Cost}(c)$. The cost-aware score applies a penalty multiplier $\lambda=0.05$:
$$
S_{\text{cost}}(c) = \frac{S(c)}{1+\lambda\widehat{\Cost}(c)}.
$$

Unused budget is dynamically tracked as $\mathbf{B}_{\text{unused}} = \mathbf{B}_{\max} - \mathbf{B}_{\text{spent}}$. Instead of halting early, the system reallocates $\mathbf{B}_{\text{unused}}$ exclusively to branches satisfying absolute high-confidence criteria:
$$
N(c) \geq 0.60, \quad RT(c) \geq 0.75, \quad \hat{S}(c) - u_S(c) \geq \tau_S.
$$

\section{Adaptive Proof Controller}

Lean failure may stem from simple syntax errors, false theorems, or missing intermediate lemmas. The system utilizes an Adaptive Proof Controller to manage the formalization ladder ($L(c)$):
$$
L(c) \in \{0.00, 0.25, 0.50, 0.70, 0.85, 1.00\}.
$$

If the compiler fails, the controller selects an action:
$$
\begin{aligned}
a_{\text{proof}}(c) \in \{ & \textsc{RepairSyntax}, \textsc{AddImports}, \textsc{SplitLemma}, \textsc{WeakenAssumptions}, \\
& \textsc{StrengthenDefinitions}, \textsc{SearchLibrary}, \textsc{DowngradeToConjecture}, \textsc{RejectFalse} \}
\end{aligned}
$$

\textbf{Lemma Decomposition Rule:} Crucially, if the Lean compiler fails repeatedly on the same theorem, the controller must attempt $\textsc{SplitLemma}$ (decomposing the proof into smaller sub-lemmas) before rejection.

\section{Bounded Auto Proof Splitting}

Formal proof failure does not necessarily imply that a mathematical claim is false. A Lean failure may arise from missing imports, monolithic theorem structure, insufficient assumptions, weak definitions, or absence of intermediate lemmas. The system therefore supports bounded proof splitting as a controlled debugging and decomposition mechanism.

For the paper-factory operating mode, proof splitting is autonomous by default but tightly bounded. Human approval is reserved for high-risk, ambiguous, or costly splits.

\subsection{Split Modes}
The proof controller supports three split modes:
$$
\operatorname{SplitMode} \in \{ \textsc{AutoBounded}, \textsc{Approval}, \textsc{Disabled} \}.
$$

The default implementation mode is:
$$
\operatorname{SplitMode}=\textsc{AutoBounded}.
$$

In $\textsc{AutoBounded}$ mode, the system may enqueue child lemmas without human approval only when structural safety, risk, and resource constraints pass. In $\textsc{Approval}$ mode, the system pauses after generating a split proposal and requires explicit human approval before child lemmas are added to the active proof tree. In $\textsc{Disabled}$ mode, failed theorem attempts may be repaired, weakened, downgraded, or rejected, but not split.

\subsection{Split Proposal}
When a theorem $T$ fails repeatedly in Lean, the Adaptive Proof Controller may propose a split:
$$
\operatorname{Split}(T)=\{T_1,\dots,T_k\},
$$
where each $T_i$ is a proposed child lemma.

A split proposal is represented as:
$$
\pi_{\text{split}}=(T_{\text{parent}},\{T_1,\dots,T_k\},E_{\text{dep}},r_{\text{split}},d_{\text{split}},h_{\text{parent}},\operatorname{SplitRisk}).
$$
Here $E_{\text{dep}}$ is the dependency edge set between child lemmas, $r_{\text{split}}$ is the reason for the split, $d_{\text{split}}$ is the current split depth, $h_{\text{parent}}$ is the ledger commit hash of the parent theorem attempt, and $\operatorname{SplitRisk}$ estimates circularity, vagueness, cost, and over-decomposition risk.

The system may propose a split only after at least one of:
$$
\{\textsc{RepairSyntax},\quad \textsc{AddImports},\quad \textsc{SearchLibrary}\}
$$
has failed to resolve the proof attempt, or the theorem is diagnosed as structurally monolithic.

\subsection{AutoBounded Split Rule}
A split may execute autonomously only if:
$$
\operatorname{AutoSplitAllowed}(\pi_{\text{split}})=1,
$$
where:
$$
\operatorname{AutoSplitAllowed}(\pi_{\text{split}})=1
\Longleftrightarrow
\begin{cases}
\operatorname{ValidSplit}(\pi_{\text{split}})=1,\\
k\leq K_{\text{split}},\\
d_{\text{split}}\leq D_{\text{split}},\\
\operatorname{CycleRisk}(\pi_{\text{split}})=0,\\
\operatorname{RepeatedSignature}(\pi_{\text{split}})=0,\\
\operatorname{SplitRisk}(\pi_{\text{split}})\leq \tau_{\text{split-risk}},\\
\operatorname{Cost}(\pi_{\text{split}})\leq B_{\text{split}}.
\end{cases}
$$

Default bounds are:
$$
K_{\text{split}}=3,
\qquad
D_{\text{split}}=2,
\qquad
\tau_{\text{split-risk}}=0.30.
$$

If $\operatorname{AutoSplitAllowed}=0$ but the theorem branch remains high-value, the branch is marked:
$$
\textsc{HumanTailEscalation}.
$$
Otherwise, the controller must autonomously choose one of:
$$
\{\textsc{WeakenAssumptions},\textsc{Simplify},\textsc{DowngradeToConjecture},\textsc{RejectFalse}\}.
$$

\subsection{Split Safety Constraints}
A split proposal is valid only if all structural safety checks pass:
$$
\operatorname{ValidSplit}(\pi_{\text{split}})=1.
$$
The validation rule is:
$$
\operatorname{ValidSplit}(\pi_{\text{split}})=1 \Longleftrightarrow
\begin{cases}
\text{no duplicate child identifiers},\\
\text{no self-dependencies},\\
\text{no unknown dependencies},\\
\text{no dependency cycles},\\
k\leq K_{\text{split}},\\
d_{\text{split}}\leq D_{\text{split}},\\
\operatorname{RepeatedSignature}(\pi_{\text{split}})=0.
\end{cases}
$$

The split dependency graph $G_{\text{split}}=(V_{\text{split}},E_{\text{dep}})$ must be acyclic, with:
$$
V_{\text{split}}=\{T_1,\dots,T_k\}.
$$

\subsection{LLM Proof Critic as Triage Only}
A read-only LLM proof critic may evaluate a split proposal before execution:
$$
\operatorname{CriticScore}(\pi_{\text{split}})\in[0,1].
$$
The critic may reduce unnecessary human escalations by flagging obvious circularity, triviality, missing dependencies, or repeated decompositions. It is not a proof authority:
$$
\text{LLM Critic}\not\Rightarrow\textsc{LeanVerified}.
$$
Only Lean compilation without forbidden tokens may create $\textsc{LeanVerified}$ status.

\subsection{Repeated Split Detection}
To avoid proof-search loops, the system computes a split signature:
$$
\operatorname{Sig}(\pi_{\text{split}})=\Hash(T_{\text{parent}},\{T_1,\dots,T_k\},E_{\text{dep}}).
$$
If the same or near-identical split signature appears more than $R_{\text{sig}}$ times, the branch is labeled:
$$
\textsc{SplitLoopDetected}.
$$
Default:
$$
R_{\text{sig}}=1.
$$

\subsection{Parent Recompilation Rule}
Successful child lemmas do not imply that the parent theorem has been proved. They only increase the feasibility of the parent proof. The parent theorem is verified only if the composed Lean file compiles:
$$
V_{\text{math\_final}}(T_{\text{parent}})=1.00
\Longleftrightarrow
\operatorname{LeanCompile}(T_{\text{parent}},T_1,\dots,T_k)=0
$$
and no forbidden token occurs:
$$
\{\texttt{sorry},\texttt{admit},\texttt{axiom},\texttt{unsafe}\}.
$$
Thus, the system may prove all child lemmas and still fail to prove the parent theorem. In that case, the parent remains labeled $\textsc{Conjecture}$ unless a complete Lean proof of the parent theorem is obtained.

\subsection{Proof Splitting Algorithm}
\begin{algorithm}[H]
\caption{Bounded Auto Proof Splitting}
\begin{algorithmic}[1]
\Require Failed theorem attempt $T$, candidate branch $c$, split mode $m$, ledger $\mathcal{G}_{\text{ledger}}$
\Ensure Updated proof state and ledger commits
\State Diagnose Lean failure for $T$
\If{failure is syntactic or import-related}
\State Execute \textsc{RepairSyntax}, \textsc{AddImports}, or \textsc{SearchLibrary}
\ElsIf{theorem is monolithic or missing intermediate lemmas}
\State Generate split proposal $\pi_{\text{split}}$
\State Validate $\operatorname{ValidSplit}(\pi_{\text{split}})$
\If{$\operatorname{ValidSplit}(\pi_{\text{split}})=0$}
\State Mark branch as $\textsc{InvalidSplit}$ or $\textsc{AutoSplitRejected}$
\ElsIf{$m=\textsc{AutoBounded}$ and $\operatorname{AutoSplitAllowed}(\pi_{\text{split}})=1$}
\State Enqueue child lemmas $\{T_1,\dots,T_k\}$ autonomously
\ElsIf{$\operatorname{HumanRequired}(c,t)=1$}
\State Mark branch as $\textsc{HumanTailEscalation}$ and emit minimal steering packet
\Else
\State Execute \textsc{WeakenAssumptions}, \textsc{DowngradeToConjecture}, or \textsc{RejectFalse}
\EndIf
\Else
\State Execute \textsc{DowngradeToConjecture} or \textsc{RejectFalse}
\EndIf
\State Commit all actions to $\mathcal{G}_{\text{ledger}}$
\end{algorithmic}
\end{algorithm}

\section{Lean Theorem Registry and Readable Lemma Wiki}

Verified mathematical artifacts should be reusable across branches. The system therefore maintains a machine-readable Lean theorem registry and a human-readable lemma wiki.

The theorem registry is the source of machine-checkable proof dependencies. The wiki is a navigation and explanation layer. The vector index is only a retrieval aid and is never treated as a source of truth.
$$
\text{Lean Registry} = \text{truth layer},
$$
$$
\text{Markdown Wiki} = \text{navigation layer},
$$
$$
\text{Vector Index} = \text{search layer},
$$
$$
\text{Ledger} = \text{provenance layer}.
$$

\subsection{Lean Theorem Record}
A theorem record is:
$$
\begin{aligned}
\tau_i = ( &\text{name}_i, \text{statement}_i, \text{namespace}_i, \text{lean\_file}_i, \text{imports}_i, \text{proof\_hash}_i, \\
&\text{commit\_hash}_i, \text{assumptions}_i, \text{lean\_version}_i, \text{mathlib\_version}_i, \text{status}_i ). 
\end{aligned}
$$

The status field is:
$$
\text{status}_i \in \{ \textsc{LeanVerified}, \textsc{Conjecture}, \textsc{Deprecated}, \textsc{Invalidated} \}.
$$

Only theorems with status $\textsc{LeanVerified}$ may be imported as dependencies:
$$
\operatorname{ImportAllowed}(\tau_i)=1 \Longleftrightarrow \tau_i.\text{status}=\textsc{LeanVerified}.
$$

\subsection{Promotion Rule}
A theorem is promoted into the registry only if:
$$
\operatorname{LeanExitCode}(\tau_i)=0
$$
and the corresponding Lean artifact contains none of:
$$
\{\texttt{sorry},\texttt{admit},\texttt{axiom},\texttt{unsafe}\}.
$$

When promoted, the system stores:
$$
\text{Lean file},\quad \text{theorem statement},\quad \text{imports},\quad \text{proof hash},\quad \text{commit hash},\quad \text{environment metadata}.
$$

\subsection{Readable Lemma Page}
For each promoted theorem $\tau_i$, the system creates a markdown lemma page:
$$
\begin{aligned}
m_i = ( & \text{title}_i, \text{informal statement}_i, \text{formal statement}_i, \text{assumptions}_i, \\
& \text{proof sketch}_i, \text{dependencies}_i, \text{origin branch}_i, \text{ledger link}_i )
\end{aligned}
$$

The default storage location is:
$$
\texttt{documents/lemmas/}.
$$

The markdown page is not a proof artifact. It is a human-readable explanation of the proof artifact. If the markdown page and Lean file disagree, the Lean file and ledger metadata dominate.

\subsection{Dependency Tracking}
If a branch $c_j$ imports theorem $\tau_i$, the ledger records:
$$
\operatorname{DependsOn}(c_j,\tau_i,h_i),
$$
where $h_i$ is the commit hash of the theorem-producing branch.

The dependency set of $c_j$ is:
$$
\operatorname{Deps}(c_j) = \{\tau_i:\operatorname{DependsOn}(c_j,\tau_i,h_i)\}.
$$

If a theorem is deprecated or invalidated, dependent branches must be rechecked:
$$
\tau_i.\text{status}\in\{\textsc{Deprecated},\textsc{Invalidated}\} \Rightarrow c_j.\text{status}\leftarrow\textsc{RecheckRequired}
$$
for all $c_j$ such that $\tau_i\in\operatorname{Deps}(c_j)$.

\subsection{Reuse Feasibility Boost}
The registry improves search by increasing proof feasibility for branches that can reuse verified lemmas. Define:
$$
\operatorname{ReuseScore}(c,\mathcal{L}_{\text{Lean}}) = \max_{\tau_i\in\mathcal{L}_{\text{Lean}}} \left( 0.40S_{\text{object}}(c,\tau_i) + 0.30S_{\text{assumption}}(c,\tau_i) + 0.30S_{\text{goal}}(c,\tau_i) \right).
$$

The proof feasibility score is updated by:
$$
F_{\text{proof}}^{+}(c) = \clip \left( F_{\text{proof}}(c) + \gamma\operatorname{ReuseScore}(c,\mathcal{L}_{\text{Lean}}), 0, 1 \right).
$$

A branch may receive a reuse boost only from theorems with status $\textsc{LeanVerified}$.

\section{Adaptive Experiment Controller}

The system treats experiments as an adaptive search tree rather than one-shot validation events. After every empirical run, the system computes an experiment state:
$$
X(c) = (\bar{r}, \operatorname{LCB}_{95}, \sigma_r, A_{\text{pass}}, B_{\text{strong}}, D_{\text{coverage}}, M_{\text{valid}}, \delta_{\text{data}}).
$$

The controller selects an action $a_{\text{exp}}(c)$ from the operational space:
$$
\begin{aligned}
a_{\text{exp}}(c) \in \{ &\textsc{Continue}, \textsc{GenerateSyntheticData}, \textsc{AddSeeds}, \textsc{AddSyntheticRegime}, \\
&\textsc{AddAblation}, \textsc{StrengthenBaseline}, \textsc{RefineHypothesis}, \\
&\textsc{DowngradeClaim}, \textsc{SpawnVariant}, \textsc{DeferRealData}, \\
&\textsc{StopSuccess}, \textsc{StopFailure} \}.
\end{aligned}
$$

The default execution policy maps the state to an action:
$$
\footnotesize
\begin{aligned}
a_{\text{exp}}(c)=
\begin{cases}
\textsc{DeferRealData}, & \delta_{\text{data}}(c)\in\{\textsc{PublicDownload},\textsc{UserProvided}\}\ \text{under MVP constraints},\\
\textsc{GenerateSyntheticData}, & \delta_{\text{data}}(c)=\textsc{SyntheticOnly}\ \text{and no generator artifact exists},\\
\textsc{StopSuccess}, & \bar{r}\geq\Delta,\ \operatorname{LCB}_{95}\geq0,\ A_{\text{pass}}=1,\ B_{\text{strong}}=1,\ M_{\text{valid}}=1,\\
\textsc{AddSeeds}, & \bar{r}\geq\Delta,\ \operatorname{LCB}_{95}<0,\\
\textsc{AddSyntheticRegime}, & \bar{r}\geq\Delta,\ D_{\text{coverage}}<\tau_D,\ \delta_{\text{data}}(c)=\textsc{SyntheticOnly},\\
\textsc{AddAblation}, & \bar{r}\geq\Delta,\ A_{\text{pass}}=0,\\
\textsc{StrengthenBaseline}, & \bar{r}\geq\Delta,\ B_{\text{strong}}=0,\\
\textsc{RefineHypothesis}, & 0<\bar{r}<\Delta,\ \operatorname{LCB}_{95}\geq0,\\
\textsc{SpawnVariant}, & \bar{r}<\Delta,\ \text{error analysis identifies localized failure},\\
\textsc{StopFailure}, & \bar{r}<0,\ \operatorname{LCB}_{95}<0,\ \text{no repairable mode exists}.
\end{cases}
\end{aligned}
$$

A final empirical claim may be written into the manuscript only after the controller reaches either $\textsc{StopSuccess}$ or $\textsc{StopFailure}$. In MVP factory mode, $\textsc{StopSuccess}$ for synthetic experiments produces $\textsc{SyntheticExperimentVerified}$, not unrestricted real-data validation.


\section{Data Requirement Regimes and MVP Data Gate}

Data access is a major determinant of autonomy. To preserve the paper-factory objective, the system explicitly classifies every candidate by its data requirement before expensive empirical validation. The data regime is:
$$
\delta_{\text{data}}(c)\in
\{\textsc{NoData},\textsc{SyntheticOnly},\textsc{PublicDownload},\textsc{UserProvided}\}.
$$

The four regimes are interpreted as follows:
\begin{center}
\begin{tabular}{p{0.24\linewidth} p{0.66\linewidth}}
\toprule
Regime & Meaning \\
\midrule
\textsc{NoData} & The paper is theoretical, conceptual, formal, or proof-based and requires no empirical dataset. \\
\textsc{SyntheticOnly} & Controlled simulations, toy examples, generated benchmarks, or synthetic counterexamples are sufficient for the intended claim. \\
\textsc{PublicDownload} & The claim requires real data, but the agent can obtain it from a public source, subject to licensing and reproducibility checks. \\
\textsc{UserProvided} & The claim requires private, proprietary, manually curated, or otherwise unavailable data that must be supplied by the user. \\
\bottomrule
\end{tabular}
\end{center}

For the initial implementation, factory mode restricts the admissible empirical scope to branches requiring at most synthetic data:
$$
\operatorname{DataAdmissible}_{\text{MVP}}(c)=1
\Longleftrightarrow
\delta_{\text{data}}(c)\in\{\textsc{NoData},\textsc{SyntheticOnly}\}.
$$

If a branch requires public or user-provided data, it is not treated as a failure of the idea. It is deferred:
$$
\delta_{\text{data}}(c)\in\{\textsc{PublicDownload},\textsc{UserProvided}\}
\Rightarrow
c.\text{status}\leftarrow\textsc{DeferredRealDataCandidate}.
$$

A user-provided data requirement is also a tail-case human escalation condition:
$$
\delta_{\text{data}}(c)=\textsc{UserProvided}
\Rightarrow
\operatorname{ExternalAccessRequired}(c,t)=1.
$$

\subsection{Synthetic Evidence Boundary}

Synthetic data may support controlled-regime, methodological, simulation, benchmark-generator, counterexample, or proof-of-concept claims. It may not support unrestricted real-world empirical claims.

The system therefore distinguishes:
$$
\textsc{SyntheticExperimentVerified}
\quad\text{from}\quad
\textsc{RealDataExperimentVerified}.
$$

Synthetic validation implies only controlled-regime evidence:
$$
\delta_{\text{data}}(c)=\textsc{SyntheticOnly}
\Rightarrow
\textsc{SyntheticExperimentVerified}(c),
$$
but:
$$
\textsc{SyntheticExperimentVerified}(c)
\not\Rightarrow
\textsc{RealDataExperimentVerified}(c).
$$

Thus, a synthetic experiment can justify claims such as ``the method succeeds under the declared generative model,'' ``the estimator recovers the target under controlled assumptions,'' or ``the failure mode appears in constructed examples.'' It cannot justify claims such as ``the method works on real markets,'' ``the model explains real mobility,'' or ``the intervention improves real-world outcomes'' unless real data validation is later added.

\subsection{Synthetic Data Artifact Requirements}

Every synthetic experiment must store a generator contract:
$$
g=(\text{generator\_code},\text{parameters},\text{random\_seeds},\text{assumptions},\text{metrics},\text{baselines},\text{hash}).
$$

The experiment artifact set must include:
$$
A_{\text{synthetic}}=(g,\text{config},\text{seeds},\text{outputs},\text{metrics},\text{plots},\text{logs}).
$$

A synthetic result is admissible only if the generator is reproducible and linked to the ledger:
$$
\begin{aligned}
\operatorname{SyntheticAdmissible}(c)=1
\Longleftrightarrow
&\operatorname{GeneratorHashValid}(c)=1 \\
&\wedge\ \operatorname{SeedsRecorded}(c)=1 \\
&\wedge\ \operatorname{MetricsPredeclared}(c)=1.
\end{aligned}
$$

\subsection{MVP Data Gate}

Before Stage C empirical validation, the system applies the data gate:
$$
\operatorname{DataGate}_{\text{MVP}}(c)=1
\Longleftrightarrow
\operatorname{DataAdmissible}_{\text{MVP}}(c)=1.
$$

If the gate fails, the branch is labeled:
$$
\textsc{RequiresRealData}
$$
or:
$$
\textsc{DeferredRealDataCandidate},
$$
and the orchestrator continues with branches that satisfy the MVP data constraint. This keeps early factory-mode execution autonomous, reproducible, and independent of user-supplied datasets.

\section{Reviewer Disagreement Resolver}

Reviewer disagreement $D_{\text{rev}}(c) = \operatorname{Var}(s_1(c),\dots,s_m(c))$ is an ambiguous signal. It may indicate paradigm-shifting novelty or fatal ambiguity. 

The Resolver explicitly classifies the variance:
$$
\operatorname{DisagreeType}(c) \in \{\textsc{NovelControversy}, \textsc{AmbiguousClaim}, \textsc{ReviewerError}, \textsc{FatalConfusion}\}.
$$

The controller enforces the following policy mappings:
\begin{itemize}[leftmargin=*]
    \item $\textsc{NovelControversy} \Rightarrow$ Preserve branch and explicitly increase epistemic uncertainty tolerance.
    \item $\textsc{AmbiguousClaim} \Rightarrow$ Route directly to the Adaptive Question Refiner.
    \item $\textsc{FatalConfusion} \Rightarrow$ Reject branch.
    \item $\textsc{ReviewerError} \Rightarrow$ Exclude the anomalous reviewer node and recompute aggregate scores.
\end{itemize}

\section{Adversarial Hardening and Evaluation Metrology}

\subsection{Adversarial Red Team Validation Layer}

To protect Stage C execution from processing conceptually hollow candidates, the Red Team Validation Layer executes four targeted semantic attacks against every Stage B survivor $c$:

\begin{enumerate}[leftmargin=*]
    \item \textbf{Novelty Attack ($RT_{\text{novelty}}$):} Targeted semantic search optimized to find counter-evidence or identical prior completions.
    \item \textbf{Bridge Attack ($RT_{\text{bridge}}$):} Operational deconstruction of cross-domain mappings. It penalizes missing mappings from primitives to objects, lack of measurable data instantiation, and decorative use of mathematical vocabulary. 
    \item \textbf{Baseline Attack ($RT_{\text{baseline}}$):} Algorithmic identification of non-trivial baselines that out-perform the proposed methodology.
    \item \textbf{Triviality Attack ($RT_{\text{triviality}}$):} Disallows theorems unless they pass a relativized significance test against baseline results and minimal examples ($T_{\text{sig}}(c) \geq 0.70$).
\end{enumerate}

The bridge survival score is defined as:
$$
B_{\text{surv}}(D,M,Q) = 0.25B_{\text{map}} + 0.20B_{\text{transfer}} + 0.20B_{\text{baseline}} + 0.15B_{\text{data}} + 0.10B_{\text{falsify}} + 0.10B_{\text{nondecorative}}.
$$

The bridge survives if $B_{\text{surv}}(D,M,Q) \geq 0.70$. We set the indicator component:
$$
RT_{\text{bridge}}(c) = \mathbf{1}\{B_{\text{surv}}(D,M,Q) \geq 0.70\}.
$$

Entry to Stage C requires $RT(c) \ge 0.75$.

\subsection{Adaptive Bridge Controller}

If a branch initially fails the Bridge Attack ($RT_{\text{bridge}}(c) = 0$), the system attempts one repair before final rejection. The bridge state is:
$$
B_{\text{state}}(c) = (B_{\text{map}}, B_{\text{transfer}}, B_{\text{baseline}}, B_{\text{data}}, B_{\text{falsify}}, B_{\text{nondecorative}}).
$$

The controller selects a reparative action:
$$
\begin{aligned}
a_{\text{bridge}}(c) \in \{ & \textsc{DefineObjects}, \textsc{ChangeMetric}, \textsc{AddBaseline}, \textsc{AddDataset}, \\
& \textsc{NarrowClaim}, \textsc{ReplaceMethod}, \textsc{RejectBridge} \}
\end{aligned}
$$

After the repair loop executes, the system recomputes $B_{\text{surv}}(D,M,Q)$. If $B_{\text{surv}}(D,M,Q) < 0.70$, the branch is formally labeled $\textsc{FalseBridge}$ and removed from the active frontier as $\textsc{RejectedRedTeam}$.

\subsection{Decision Metrology and Score Uncertainty Bounding}

Candidate assessments are treated as continuous distributions with a mean expected score $\hat{S}(c)$ and epistemic uncertainty $u_S(c)$:
$$
S(c) = \hat{S}(c) \pm u_S(c).
$$

To ensure survival, a candidate must pass the conservative lower-bound threshold:
$$
\hat{S}(c) - u_S(c) \geq \tau_S.
$$
If this inequality fails, the candidate is marked as \textsc{PrunedUncertain} in the ledger and removed from the active search frontier.

\subsection{Parametric Calibration and MVP Priors}

All numerical thresholds and heuristic weights presented in this specification are \textit{Minimum Viable Product (MVP) priors}. They are reasonable starting boundaries designed to make the architecture concrete and testable, not fixed theoretical constants. 

The numerical values fall into three distinct classes:
\begin{itemize}[leftmargin=*]
    \item \textbf{Hard Correctness Gates:} (e.g., $V_{\text{math\_final}} = 1.00$, no forbidden tokens). Very important; these absolutely protect the integrity of empirical and mathematical truth claims.
    \item \textbf{Search-Control Thresholds:} (e.g., $N(c) \geq 0.45$, $F(c) \geq 0.80$, $RT(c) \geq 0.75$). Important, but tunable. They govern branch pruning and exploration versus exploitation.
    \item \textbf{Heuristic Weights:} (e.g., $0.25$ novelty, $0.20$ feasibility). Least fixed; these balance the scoring components.
\end{itemize}

Let the tunable search-control thresholds be defined as the parameter vector:
$$
\theta = (\tau_{\text{opportunity}}, \tau_{\text{novelty}}, \tau_{\text{feasibility}}, \tau_{\text{verifiability}}, \tau_{\text{redteam}}, \tau_{\text{adequacy}}, \tau_S).
$$

The core philosophy of the system is: \textit{The architecture fixes the structure of the decision process. The calibration suite fixes the numerical thresholds.}

To systematically optimize these activation thresholds, the multi-agent architecture is continuously evaluated against a Calibration Benchmark Suite $\mathcal{D}_{\text{cal}}$. This reference set is uniformly partitioned across four true-label categories:
$$
\mathcal{D}_{\text{cal}} = \mathcal{D}_{\text{success}} \cup \mathcal{D}_{\text{forced}} \cup \mathcal{D}_{\text{known}} \cup \mathcal{D}_{\text{impossible}}.
$$
Where each subset contains $N=20$ curated structural problems:
\begin{itemize}[leftmargin=*]
    \item $\mathcal{D}_{\text{success}}$: Documented, highly valid cross-domain mathematical integrations.
    \item $\mathcal{D}_{\text{forced}}$: Conceptually weak, forced semantic analogies designed to mimic valid literature.
    \item $\mathcal{D}_{\text{known}}$: Existing scientific breakthroughs used to test over-claiming thresholds.
    \item $\mathcal{D}_{\text{impossible}}$: Conceptually flawed or structurally unprovable research hypotheses.
\end{itemize}

The system systematically learns the optimal threshold configuration $\theta^\star$ by minimizing the empirical classification error over the calibration benchmark $\mathcal{D}_{\text{cal}}$:
$$
\theta^\star = \arg\min_{\theta} \operatorname{Error}(\theta;\mathcal{D}_{\text{cal}}).
$$




\section{Strategic Questioner Controller}

The framework already contains local controllers for question refinement, literature retrieval, bridge repair, proof search, experiment design, disagreement resolution, and abstraction. However, local controllers alone are not sufficient to reproduce the scientific judgment required to turn many explored branches into one meaningful research output. The system must explicitly ask whether the current branch, paper nucleus, or surviving branch portfolio is clear, novel, repairable, sufficiently evidenced, simple enough, and worth continuing.

The Strategic Questioner Controller is a lightweight meta-controller. It is not a free-form brainstorming agent, not a publication-logistics module, and not a frequent human-question generator. In particular, questions about venue selection, journal fit, formatting style, or submission packaging are outside the scientific discovery core. Such questions may be handled by an optional later publication layer, but they may not influence verification labels, evidence requirements, or scientific truth claims.

The controller converts explicit diagnostic questions into routing actions for existing controllers. In factory mode, its default behavior is autonomous routing: questions are asked internally, converted into actions, and committed to the ledger. It routes to 	extsc{AskHuman} only when the autonomy contract says that human intervention is required. Its purpose is to prevent three failure modes:
\begin{enumerate}[leftmargin=*]
    \item \textbf{Premature convergence:} selecting the highest-scoring local branch before asking whether a stronger formulation, repair, or abstraction exists.
    \item \textbf{Unbounded elaboration:} adding more methods, lemmas, data, or experiments when the branch should instead be simplified, downgraded, or stopped.
    \item \textbf{Weak synthesis:} producing a narrow paper from one branch when several explored branches reveal a more meaningful general model.
\end{enumerate}

The controller is invoked through a staged question-selection policy. The full question bank is \emph{not} asked at every step. Instead, the system performs cheap invariant checks after every state-changing action, stage-specific strategic questioning at major gates, and specialized questioning only when diagnostic triggers fire.

\subsection{Question Invocation Policy}

Let $s_t$ be the current run state, $\operatorname{Stage}(t)$ the current pipeline stage, and $\operatorname{Triggers}(s_t)$ the set of active diagnostic triggers. The question set evaluated at time $t$ is:
$$
Q_t = \operatorname{SelectQuestions}\bigl(s_t,\operatorname{Stage}(t),\operatorname{Triggers}(s_t)\bigr).
$$
The resulting answers are routed to controller actions:
$$
a_{\mathrm{meta}}(t)=\operatorname{RouteAnswer}(Q_t,c_t,P_t,\mathcal{B}_t).
$$
Thus, the framework asks questions continuously, but it does not ask all questions continuously. Questioning is budgeted and context-dependent:
$$
Q_t = Q_{\mathrm{micro}}(t)\cup Q_{\mathrm{stage}}(t)\cup Q_{\mathrm{trigger}}(t).
$$

\subsubsection*{Micro-checks after every state-changing action}
After every controller action, model call, artifact creation, branch mutation, proof attempt, experiment run, human command, or scoring update, the system asks only the invariant micro-checks:
\begin{enumerate}[leftmargin=*]
    \item Does the branch still satisfy the user constraint set $C$?
    \item Did the action improve the branch, reduce uncertainty, or produce useful evidence?
    \item Did the action add unnecessary complexity without improving novelty, feasibility, verifiability, or meaning?
    \item Was the action recorded as a ledger commit with reason code and artifact references?
    \item Were all newly produced artifacts hashed and linked to the producing commit?
    \item Did any verification label change, and if so, is the new label justified by evidence rather than prose?
    \item Did the action create a new blocking issue requiring repair, downgrade, pruning, or human steering?
\end{enumerate}
These micro-checks protect invariants. They do not perform full scientific strategy review.

\subsubsection*{Stage-level strategic questioning}
At major stage gates, the controller selects the relevant scientific questions from the mandatory bank:

\begin{center}
\begin{tabular}{p{0.23\linewidth} p{0.67\linewidth}}
\toprule
Checkpoint & Required question families \\
\midrule
After Stage A exploration & Core contribution, clarity, novelty, scope, diversity, and early stopping questions. \\
After literature retrieval & Retrieval adequacy, terminology expansion, closest-prior comparison, novelty-risk, and retrieval-sufficiency questions. \\
After Stage B structural validation & Mathematical or empirical validity, repair sufficiency, baseline strength, bridge validity, simplicity, and paper-shape questions. \\
After red-team filtering & Fatal-flaw, repairability, claim weakening, negative-result, limitation, and evidence-boundary questions. \\
After Stage C verification & Final-label sufficiency, Lean or experiment status, missing evidence, data richness, baseline adequacy, ablation, and downgrade questions. \\
Before Abstract Synthesis & Shared-structure, instantiation-map, abstraction-compression, special-case, and global-model questions. \\
Before final manuscript generation & Final-nucleus, unsupported-claim, appendix/supporting-case, limitation, and failed-branch-retention questions. \\
\bottomrule
\end{tabular}
\end{center}

\subsubsection*{Triggered specialized questions}
Specialized questions fire only when their trigger appears:
\begin{center}
\begin{tabular}{p{0.27\linewidth} p{0.58\linewidth}}
\toprule
Trigger & Specialized questions \\
\midrule
High complexity or many objects & Is this overcomplicated? Can the claim be stated with fewer objects, assumptions, parameters, lemmas, datasets, or metrics? \\
Weak or narrow data & Is richer data needed? Should additional datasets, regimes, seeds, time periods, samples, or robustness tests be added? \\
Weak baseline or easy counterexample & Is the baseline strong enough? What is the strongest simple counterexample, prior method, or naive baseline? \\
Recent repair applied & Did the repair actually fix the stated flaw? Did it introduce new assumptions or merely hide the problem? \\
Lean failure & Is the failure syntactic, definitional, structural, caused by a false theorem, or caused by missing intermediate lemmas? \\
Experiment instability & Are more seeds, ablations, stronger baselines, richer data, or claim downgrades required? \\
Many related surviving branches & Can these branches be unified by a global abstraction $G$ with explicit instantiation maps $\psi_i:G\to c_i$? \\
Many redundant branches & Should branches be merged, clustered, or reduced to one representative branch plus supporting cases? \\
Repeated repair or proof loop & Is more budget likely to change the conclusion, or should the branch be downgraded, stopped, or escalated to the human operator? \\
Pre-synthesis uncertainty & Should the final output be centered on a global abstraction $G^\star$, one winning branch $c^\star$, a negative result, or a failed-search report? \\
\bottomrule
\end{tabular}
\end{center}

\subsubsection*{Question budget}
To prevent reflective overhead, the controller may impose a maximum number of strategic questions per invocation:
$$
|Q_t\setminus Q_{\mathrm{micro}}(t)|\leq q_{\max}.
$$
The default policy prioritizes failed diagnostics, recently changed fields, and questions whose answers can change the next controller action. A question is not asked if its answer cannot affect routing, scoring, verification status, synthesis, or stopping.

\subsubsection*{Interpretation}
The system therefore uses questions as a control mechanism rather than as open-ended discussion. The invariant micro-checks run after every state-changing step. The stage-level questions run at gates. The specialized questions run only when triggered. This preserves the benefit of scientific self-questioning without turning the framework into an expensive loop of generic reflection.

\subsection{Three Levels of Scientific Questioning}

The controller operates at three levels:
$$
Q_{\mathrm{strategic}}=(Q_{\mathrm{branch}},Q_{\mathrm{paper}},Q_{\mathrm{portfolio}}).
$$

\begin{itemize}[leftmargin=*]
    \item $Q_{\mathrm{branch}}$ asks whether an individual candidate branch is valid, novel, repairable, evidenced, and worth continuing.
    \item $Q_{\mathrm{paper}}$ asks whether a manuscript nucleus has the right core contribution, scope, framing, and synthesis structure.
    \item $Q_{\mathrm{portfolio}}$ asks how multiple surviving branches should be merged, split, abstracted, deprioritized, or converted into separate outputs.
\end{itemize}

For a candidate branch $c$, manuscript nucleus $P$, and surviving branch set $\mathcal{B}$, the diagnostic state is:
$$
\begin{aligned}
Q_{\mathrm{meta}}(c,P,\mathcal{B})=(
&C_{\mathrm{core}},C_{\mathrm{clarity}},C_{\mathrm{novelty}},C_{\mathrm{valid}},C_{\mathrm{repair}},C_{\mathrm{suff}},\\
&C_{\mathrm{simp}},C_{\mathrm{data}},C_{\mathrm{base}},C_{\mathrm{scope}},C_{\mathrm{meaning}},C_{\mathrm{abs}},C_{\mathrm{portfolio}},C_{\mathrm{stop}}),
\end{aligned}
$$
where each component lies in $[0,1]$. Higher values mean that the branch, paper nucleus, or portfolio satisfies the corresponding diagnostic dimension.

The output is a structured report:
$$
\operatorname{QuestionerReport}(c,P,\mathcal{B})=
(Q_{\mathrm{asked}},Q_{\mathrm{failed}},a_{\mathrm{meta}},r_{\mathrm{meta}},h_{\mathrm{commit}}),
$$
where $Q_{\mathrm{asked}}$ is the set of diagnostic questions evaluated, $Q_{\mathrm{failed}}$ is the subset whose answers are insufficient, $a_{\mathrm{meta}}$ is the selected routing action, $r_{\mathrm{meta}}$ is the reason, and $h_{\mathrm{commit}}$ is the ledger commit hash.

\subsection{Mandatory Diagnostic Question Bank}

The controller maintains the following mandatory question bank. The system does not ask every question at every invocation. Instead, it selects the relevant subset using $\operatorname{SelectQuestions}$. A question from this bank may be skipped only when its trigger is absent or when its answer cannot affect routing, scoring, verification status, synthesis, or stopping. Venue-selection and submission-logistics questions are intentionally excluded.

\subsubsection*{Core Contribution Questions}
\begin{enumerate}[leftmargin=*]
    \item What is the central contribution of this branch or paper nucleus?
    \item Is the original motivating idea fully used, or is the paper exploiting only a weaker fragment of it?
    \item Is the contribution a theorem, model, method, empirical finding, dataset, negative result, measurement principle, or unifying framework?
    \item Which claim would remain if all secondary material were removed?
    \item Does every major section support the central contribution?
    \item Is there a stronger or cleaner version of the same idea that the current branch has not yet expressed?
    \item Is the contribution genuinely scientific, or mostly a change of exposition or notation?
\end{enumerate}

\subsubsection*{Clarity and Formulation Questions}
\begin{enumerate}[leftmargin=*]
    \item Can the central claim be written in one falsifiable or verifiable sentence?
    \item What are the mathematical, empirical, or conceptual objects involved?
    \item What are the inputs, outputs, assumptions, and claimed conclusion?
    \item Which terms remain undefined, ambiguous, overloaded, or metaphorical?
    \item Is the scope narrow enough to verify within the available budget?
    \item Is the branch asking one question, or mixing several independent questions?
    \item What is the weakest part of the current formulation?
    \item Would simplifying the formulation preserve the main contribution?
\end{enumerate}

\subsubsection*{Novelty and Prior-Work Questions}
\begin{enumerate}[leftmargin=*]
    \item Is retrieval adequacy high enough for the claimed novelty level?
    \item Have adjacent fields, alternative terminology, and older literatures been searched?
    \item Could the proposed result be a rediscovery under different notation?
    \item Is the novelty in the object, method, theorem, experiment, combination, or interpretation?
    \item Is low literature density caused by true underuse or by bad retrieval?
    \item Which prior work is closest, and what exactly remains new after comparison?
    \item Does the branch need a narrower novelty claim?
    \item Should retrieval be expanded before the branch continues?
\end{enumerate}

\subsubsection*{Mathematical Validity Questions}
\begin{enumerate}[leftmargin=*]
    \item Is the formal statement mathematically well-posed?
    \item Are the assumptions sufficient for the conclusion?
    \item Are the assumptions so strong that the result becomes trivial?
    \item Does the proof strategy actually use the main idea of the paper?
    \item Is there a counterexample, degenerate case, or limiting case that defeats the claim?
    \item Can the mathematical claim be formalized in Lean without changing its meaning?
    \item If Lean fails, is the failure syntactic, definitional, structural, or evidence that the claim is false?
    \item Should the theorem be split, weakened, repaired, or downgraded to a conjecture?
\end{enumerate}

\subsubsection*{Repair Sufficiency Questions}
\begin{enumerate}[leftmargin=*]
    \item What exact flaw was the repair supposed to fix?
    \item Does the proposed repair eliminate the flaw, or only hide it behind stronger wording?
    \item Does the repaired branch still satisfy the original user constraints and scientific motivation?
    \item Did the repair introduce new assumptions, dependencies, or failure modes?
    \item Is the repaired claim weaker, stronger, or merely different from the original claim?
    \item Does the repaired proof, experiment, or argument pass the relevant gate?
    \item Should the branch undergo another repair cycle, be downgraded, or be stopped?
\end{enumerate}

\subsubsection*{Evidence and Sufficiency Questions}
\begin{enumerate}[leftmargin=*]
    \item What final label does the current evidence support: \textsc{LeanVerified}, \textsc{SyntheticExperimentVerified}, \textsc{RealDataExperimentVerified}, \textsc{Conjecture}, \textsc{NegativeResult}, \textsc{Limitation}, or \textsc{Unsupported}?
    \item What minimum missing evidence prevents promotion to a stronger label?
    \item Does the evidence support the central claim, or only a weaker subclaim?
    \item Are all supporting artifacts linked by content hash and ledger commit?
    \item Is there a predeclared proof target or experiment contract?
    \item Has the branch passed the required stage gate for its current claim strength?
    \item Would a skeptical reviewer accept the current evidence as sufficient for the stated label?
    \item Is any unsupported claim being smuggled into the narrative through prose?
\end{enumerate}

\subsubsection*{Simplicity and Overcomplication Questions}
\begin{enumerate}[leftmargin=*]
    \item Can the claim be stated with fewer objects, assumptions, or parameters?
    \item Is the chosen method necessary, or is it decorative vocabulary?
    \item Is there a simpler theorem, baseline, model, or experiment that explains the same phenomenon?
    \item Are any assumptions introduced only to make the result work?
    \item Are there redundant definitions, lemmas, datasets, metrics, or branches?
    \item Does the current abstraction compress the paper, or does it add vague complexity?
    \item What could be removed without weakening the main contribution?
    \item Should the branch be simplified before receiving more budget?
\end{enumerate}

\subsubsection*{Data, Experiment, and Baseline Questions}
\begin{enumerate}[leftmargin=*]
    \item Which data regime does the branch require: \textsc{NoData}, \textsc{SyntheticOnly}, \textsc{PublicDownload}, or \textsc{UserProvided}?
    \item Can the intended claim be verified with no data or synthetic data only?
    \item If synthetic data is sufficient, what generator, assumptions, seeds, metrics, and baselines must be declared before execution?
    \item Does the synthetic generator match the assumptions of the claim, or does it make the experiment circular?
    \item Is the branch accidentally making a real-world empirical claim from synthetic evidence only?
    \item If real data is required, should the branch be deferred as \textsc{DeferredRealDataCandidate} under the MVP policy?
    \item Are the metrics aligned with the scientific claim rather than merely convenient?
    \item What is the strongest simple baseline?
    \item What is the strongest known prior method, theorem, or empirical benchmark?
    \item Does the result survive ablations, alternative metrics, reasonable perturbations, and stronger baselines?
    \item Is there a risk of leakage, selection bias, survivorship bias, cherry-picking, or data snooping?
    \item Would adding richer data change the claim label, or only increase cost outside the MVP data boundary?
    \item Should the experiment be continued, redesigned, converted to synthetic validation, downgraded, deferred, or stopped?
\end{enumerate}

\subsubsection*{Scope, Merge, Split, and Output-Form Questions}
\begin{enumerate}[leftmargin=*]
    \item Is this material part of the main paper, an appendix, a separate paper, a lemma, an experiment, or a failed branch?
    \item Is the current manuscript too broad for one coherent contribution?
    \item Should two branches be merged because they share the same mechanism?
    \item Should a broad branch be split into simpler independently verifiable branches?
    \item Does an appendix strengthen the paper, or does it distract from the main contribution?
    \item Should an extension remain optional until the main result is stable?
    \item Is a negative result informative enough to preserve as an output?
    \item What is the smallest coherent paper that preserves the strongest contribution?
\end{enumerate}

\subsubsection*{Meaningfulness and Research-Value Questions}
\begin{enumerate}[leftmargin=*]
    \item Why does this branch matter scientifically?
    \item What problem, concept, measurement, theorem, or empirical uncertainty does it clarify?
    \item Does the branch change understanding, improve measurement, unify cases, or expose a useful limitation?
    \item Is the result nontrivial relative to the baseline and minimal examples?
    \item Is the claim strong enough to justify a paper-level output, or only a note, appendix, lemma, or failed branch?
    \item Does the branch produce reusable concepts, tools, data, proofs, or experimental protocols?
    \item Is continuing this branch more valuable than reallocating budget to another branch?
\end{enumerate}

\subsubsection*{Abstraction and Global-Synthesis Questions}
\begin{enumerate}[leftmargin=*]
    \item Do multiple surviving branches share objects, assumptions, mechanisms, metrics, or proof structure?
    \item Can a general object $G$ be stated so that the branches become true instances?
    \item For each covered branch $c_i$, is there an explicit instantiation map $\psi_i:G\rightarrow c_i$?
    \item Are the branches genuine special cases, or merely analogies?
    \item Does the abstraction reduce the complexity of the final paper?
    \item Does the abstraction generate new corollaries, experiments, predictions, or boundary cases?
    \item Does any verified result actually depend on the abstraction?
    \item Should the final manuscript be centered on a global model $G^\star$ or on a single branch $c^\star$?
\end{enumerate}

\subsubsection*{Portfolio and Prioritization Questions}
\begin{enumerate}[leftmargin=*]
    \item Among the surviving branches, which has the best novelty-to-evidence-to-effort ratio?
    \item Which branch is strongest after repairs, red-team attacks, and verification attempts?
    \item Which branch should receive the next unit of budget?
    \item Which branch should be deprioritized because it is weaker, redundant, too broad, or too costly?
    \item Does one branch provide the main paper while another becomes a supporting case?
    \item Are several branches competing formulations of the same idea?
    \item What is the next concrete research action across the whole branch portfolio?
\end{enumerate}

\subsubsection*{Stopping, Downgrading, and Human-Steering Questions}
\begin{enumerate}[leftmargin=*]
    \item What is the most important unresolved issue blocking the branch or paper nucleus?
    \item Is more budget likely to change the conclusion?
    \item Is the branch stuck in a repeated repair, retrieval, proof-splitting, or experiment loop?
    \item Should the branch continue, be simplified, be downgraded, become a negative result, or be pruned?
    \item Does the next decision require scientific taste, strategic preference, or external knowledge beyond the current evidence?
    \item Should the system ask the human to approve a proof split, abstraction, claim downgrade, budget extension, branch merge, branch split, or final nucleus?
    \item What is the minimal question that should be asked of the human operator?
\end{enumerate}

\subsection{Meta-Actions and Routing}

The controller chooses a single primary action, with optional secondary actions:
$$
\begin{aligned}
a_{\mathrm{meta}}(c,P,\mathcal{B})\in\{&\textsc{ExploitCoreIdea},\textsc{ImproveQuestion},\textsc{Simplify},\textsc{ValidateRepair},\\
&\textsc{AddSyntheticData},\textsc{DeferRealData},\textsc{StrengthenBaseline},\textsc{ImproveRetrieval},\textsc{RunAblation},\\
&\textsc{TightenProof},\textsc{AttemptAbstraction},\textsc{MergeBranches},\textsc{SplitBranch},\\
&\textsc{PromoteToMainResult},\textsc{MoveToAppendix},\textsc{DeprioritizeBranch},\\
&\textsc{DowngradeClaim},\textsc{StopSuccess},\textsc{StopFailure},\textsc{AskHuman}\}.
\end{aligned}
$$

Routing is defined by:
\begin{center}
\begin{tabular}{p{0.29\linewidth} p{0.60\linewidth}}
\toprule
Meta-action & Routed controller or effect \\
\midrule
\textsc{ExploitCoreIdea} & Route to Question Refiner or Abstract Synthesis Controller with a requirement to restate and use the strongest motivating idea. \\
\textsc{ImproveQuestion} & Route to Adaptive Question Refiner. \\
\textsc{Simplify} & Route to Question Refiner, Proof Controller, Experiment Controller, or Abstract Synthesis Controller with a simplicity constraint. \\
\textsc{ValidateRepair} & Route to the controller that produced the repair and require a before/after flaw check. \\
\textsc{AddSyntheticData} & Route to Adaptive Experiment Controller with a synthetic generator, seeds, metrics, and baselines. \\
\textsc{DeferRealData} & Label the branch as \textsc{DeferredRealDataCandidate} or \textsc{RequiresRealData} when public or user-provided data is required under MVP constraints. \\
\textsc{StrengthenBaseline} & Route to Baseline Attack, Bridge Controller, or Experiment Controller. \\
\textsc{ImproveRetrieval} & Route to Adaptive Literature Controller. \\
\textsc{RunAblation} & Route to Adaptive Experiment Controller. \\
\textsc{TightenProof} & Route to Adaptive Proof Controller or Lean Theorem Registry. \\
\textsc{AttemptAbstraction} & Route to Abstract Synthesis Controller. \\
\textsc{MergeBranches} & Merge compatible branches under a shared paper nucleus or abstraction proposal. \\
\textsc{SplitBranch} & Split an over-broad branch into independently verifiable branches. \\
\textsc{PromoteToMainResult} & Mark branch as candidate final nucleus subject to verification and synthesis gates. \\
\textsc{MoveToAppendix} & Preserve material as supporting evidence, extension, or secondary result. \\
\textsc{DeprioritizeBranch} & Reduce budget priority while preserving ledgered evidence and failure reasons. \\
\textsc{DowngradeClaim} & Change claim label to a weaker admissible label and commit the downgrade. \\
\textsc{StopSuccess} & Preserve branch as final-eligible under its verified or labeled status. \\
\textsc{StopFailure} & Mark branch as failed, negative, or limitation with reason code. \\
\textsc{AskHuman} & Emit a minimal steering packet only when $\operatorname{HumanRequired}(c,t)=1$; otherwise choose the closest autonomous action. \\
\bottomrule
\end{tabular}
\end{center}

\subsection{Decision Policy}

The default decision policy is:
$$
\footnotesize
\begin{aligned}
a_{\mathrm{meta}}(c,P,\mathcal{B})=
\begin{cases}
\textsc{ExploitCoreIdea}, & C_{\mathrm{core}}<0.70,\\
\textsc{ImproveQuestion}, & C_{\mathrm{clarity}}<0.70,\\
\textsc{ImproveRetrieval}, & C_{\mathrm{novelty}}<0.80,\\
\textsc{ValidateRepair}, & C_{\mathrm{repair}}<0.75\ \text{and a repair was recently applied},\\
\textsc{Simplify}, & C_{\mathrm{simp}}<0.60\ \text{or }C_{\mathrm{scope}}<0.60,\\
\textsc{StrengthenBaseline}, & C_{\mathrm{base}}<0.70,\\
\textsc{AddSyntheticData}, & C_{\mathrm{data}}<0.70\ \text{and }\delta_{\text{data}}(c)\in\{\textsc{NoData},\textsc{SyntheticOnly}\},\\
\textsc{DeferRealData}, & \delta_{\text{data}}(c)\in\{\textsc{PublicDownload},\textsc{UserProvided}\}\ \text{under MVP constraints},\\
\textsc{TightenProof}, & C_{\mathrm{valid}}<0.70\ \text{and }c\text{ is mathematical},\\
\textsc{AttemptAbstraction}, & C_{\mathrm{abs}}\geq0.75\ \text{and }|\mathcal{B}|\geq2,\\
\textsc{MergeBranches}, & C_{\mathrm{portfolio}}<0.70\ \text{and branches are redundant or structurally shared},\\
\textsc{SplitBranch}, & C_{\mathrm{scope}}<0.60\ \text{and the branch contains separable claims},\\
\textsc{DowngradeClaim}, & C_{\mathrm{suff}}<0.70\ \text{but a weaker claim is supported},\\
\textsc{StopFailure}, & C_{\mathrm{meaning}}<0.60\ \text{or }C_{\mathrm{stop}}<0.50,\\
\textsc{AskHuman}, & \operatorname{HumanRequired}(c,t)=1,\\
\textsc{StopSuccess}, & \text{all mandatory gates for the current label pass}.
\end{cases}
\end{aligned}
$$

Every invocation produces a ledger commit:
$$
\operatorname{QuestionerReport}(c,P,\mathcal{B})\Rightarrow \exists n_i\in\mathcal{N}_{\mathrm{ledger}}.
$$

The Strategic Questioner Controller may route, simplify, merge, split, downgrade, or stop a branch autonomously. It may pause only when $\operatorname{HumanRequired}(c,t)=1$. It may not itself verify a theorem or validate an empirical claim:
$$
\textsc{QuestionerReport}\not\Rightarrow\textsc{LeanVerified},
$$
$$
\textsc{QuestionerReport}\not\Rightarrow\textsc{ExperimentVerified}.
$$

\subsection{Excluded Publication-Packaging Questions}

The following questions are useful for a later publication layer but are excluded from the scientific discovery core:
\begin{itemize}[leftmargin=*]
    \item Which journal, conference, or venue should receive the manuscript?
    \item How should the paper be formatted for a specific venue?
    \item Which submission deadline or publication strategy should be targeted?
    \item How should the abstract or introduction be stylistically adapted for a particular venue?
\end{itemize}

A publication-packaging layer may modify exposition, ordering, emphasis, and formatting. It may not modify evidence labels, proof status, experiment status, ledger records, or scientific claims.


\section{Autonomy Contract and Tail-Case Human Escalation}

The system is intended to operate as an autonomous paper factory rather than as a step-by-step human co-pilot. Human intervention is therefore treated as a tail-case escalation, not as the default resolution mechanism.

The autonomy contract is:
$$
\operatorname{HumanRequired}(c,t)=1
$$
only if at least one of the following conditions holds:
$$
\operatorname{DecisionUncertainty}(c,t)>\tau_{\text{human}},
$$
$$
\operatorname{ActionRisk}(c,t)>\tau_{\text{risk}},
$$
$$
\operatorname{ExtraBudgetRequired}(c,t)=1,
$$
$$
\operatorname{IrreversibleDecision}(c,t)=1,
$$
$$
\operatorname{ExternalAccessRequired}(c,t)=1,
$$
or:
$$
\operatorname{UserPreferenceNeeded}(c,t)=1.
$$
Otherwise:
$$
\operatorname{HumanRequired}(c,t)=0.
$$

When $\operatorname{HumanRequired}(c,t)=0$, the orchestrator must select an autonomous action from the allowed controller action space. Ordinary research-management decisions are not escalated to the human. These include:
\begin{itemize}[leftmargin=*]
    \item narrowing, reframing, or simplifying a question;
    \item increasing retrieval adequacy;
    \item adding synthetic data, seeds, ablations, or stronger baselines;
    \item weakening assumptions or downgrading unsupported claims;
    \item stopping weak branches, converting failures to limitations, or producing negative results;
    \item selecting a final abstraction or branch when the score margin is clear.
\end{itemize}

Human escalation is allowed for tail cases:
\begin{itemize}[leftmargin=*]
    \item high-risk proof splits or abstraction choices;
    \item unresolved high-value branches with high decision uncertainty;
    \item extra compute, paid access, public-data access outside MVP scope, private data, or external tool access;
    \item irreversible authorial decisions such as final release approval;
    \item genuine ties between final nuclei where user preference is required.
\end{itemize}

The default policy is therefore:
$$
Q_t \longrightarrow a_t,
$$
not:
$$
Q_t \longrightarrow \text{ask human}.
$$
The human operator receives only a minimal steering packet:
$$
H_{\text{packet}}(c)=(q_{\text{ask}},r_{\text{block}},A_{\text{evidence}},a_{\text{suggested}},h_{\text{latest}}),
$$
where $q_{\text{ask}}$ is the single smallest question needed to unblock the tail case.

\section{Operational Risks and Mitigations}

The architecture provides disciplined exploration, labeling, provenance, and gated verification only under correct implementation of its calibration, retrieval, ledger, and controller-loop safeguards. It does not guarantee novelty, publishability, or scientific value by architecture alone.

\subsection{Calibration Sensitivity and Active Calibration Queue}
The threshold vector:
$$
\theta=(\tau_{\text{opportunity}},\tau_{\text{novelty}},\tau_{\text{feasibility}},\tau_{\text{verifiability}},\tau_{\text{redteam}},\tau_{\text{adequacy}},\tau_S,\tau_A)
$$
contains MVP priors rather than theoretical constants. To reduce overfitting to a static benchmark, the system maintains an active calibration queue.

Branches are flagged for calibration review when:
\begin{itemize}[leftmargin=*]
    \item $\texttt{human\_override}$ is used;
    \item reviewer disagreement is high;
    \item a red-team decision is later reversed;
    \item a high-scoring branch later fails;
    \item a nearly pruned branch later succeeds;
    \item an abstraction is accepted or rejected controversially.
\end{itemize}

The benchmark update rule is curated, not automatic:
$$
\mathcal{D}_{\text{cal}}^{+}
=
\mathcal{D}_{\text{cal}}
\cup
\operatorname{Reviewed}(\mathcal{D}_{\text{flagged}}).
$$
A flagged branch may influence future calibration only after review and labeling.

\subsection{Literature Blindspots}
The retrieval adequacy certificate mitigates, but does not eliminate, literature blindspots. The system therefore treats novelty as uncertain whenever semantic, keyword, citation, diversity, or adversarial retrieval channels disagree. Strong novelty claims are blocked if:
$$
\rho_{\text{adequacy}}(c)<\tau_{\text{adequacy}}.
$$
The branch is labeled $\textsc{InsufficientRetrievalAdequacy}$ rather than promoted to deep verification.

\subsection{Global Stagnation Index}
Agentic repair loops can waste budget without improving the branch. The system tracks a global stagnation index:
$$
\operatorname{Stag}(c,t)
=
\#\{\text{recent actions with }\Delta S(c)<\epsilon_S
\text{ and no verification improvement}\}.
$$

If:
$$
\operatorname{Stag}(c,t)\geq N_{\text{stag}},
$$
the branch must be resolved autonomously by one of:
$$
\{\textsc{Simplify},\textsc{DowngradeClaim},\textsc{ConvertToNegativeResult},\textsc{StopFailure}\}.
$$
Human escalation is allowed only when:
$$
S(c)>\tau_{\text{high-value}}
\quad\text{and}\quad
u_S(c)>\tau_{\text{uncertain}}.
$$

\subsection{Proof Bottlenecks}
Bounded auto proof splitting reduces human bottlenecks while preserving verification safety. Auto-splitting may enqueue child lemmas only under strict depth, width, acyclicity, signature, risk, and cost constraints. An LLM proof critic may triage a split proposal, but may not create proof authority:
$$
\text{LLM Critic}\not\Rightarrow\textsc{LeanVerified}.
$$

\subsection{Fake Abstraction Risk}
The Abstract Synthesis Controller may generate overly broad or decorative unifications. The Abstraction Attack therefore tests whether branches are genuine instances of a proposed model, whether assumptions are compatible, whether compression is real, and whether verified results depend on the abstraction. If:
$$
A(G,\mathcal{B})<\tau_A
\quad\text{or}\quad
RT_{\text{abstract}}(G)<\tau_{\text{abstract}},
$$
the abstraction is rejected and the system falls back to the best verified branch.

\subsection{Runtime Context Compression}
The ledger is never pruned. It remains the full provenance record:
$$
\mathcal{G}_{\text{ledger}}=\text{full provenance}.
$$
However, runtime context may be compressed to prevent context bloat:
$$
\operatorname{RuntimeSummary}(c)
=
\operatorname{Compress}(\operatorname{LedgerHistory}(c)).
$$
Runtime summaries are convenience artifacts only:
$$
\operatorname{RuntimeSummary}(c)\not\Rightarrow\operatorname{Provenance}(c).
$$
Only the full ledger and artifact hashes define provenance.

\subsection{Operational Principle}
The operational principle is:
$$
\boxed{\text{Autonomous by default; human intervention only for tail-risk, high-ambiguity, high-cost, or irreversible decisions.}}
$$
Operational mitigations must reduce loops, blindspots, and bottlenecks without weakening verification authority.


\section{Abstract Synthesis Controller}

The preceding controllers operate primarily at the branch level. They refine questions, retrieve literature, repair bridges, verify proofs, and validate experiments for individual candidates. However, a strong final paper often requires a higher-level synthesis step: the system must determine whether several surviving branches are instances of a single general model rather than merely selecting the highest-scoring branch.

The Abstract Synthesis Controller therefore operates after Stage C verification and before final manuscript synthesis. Its purpose is to compress the surviving research tree into the most meaningful paper-level nucleus.

Let the surviving branch set be:
$$
\mathcal{B}=\{c_1,\dots,c_k\}.
$$

Each branch contains structured objects, assumptions, methods, claims, results, and verification evidence:
$$
c_i=(O_i,A_i,M_i,Q_i,H_i,R_i,V_i).
$$

The controller proposes a candidate abstraction:
$$
G=(O_G,A_G,M_G,\Theta_G,\mathcal{I}_G),
$$
where $O_G$ are general objects, $A_G$ are general assumptions, $M_G$ is the general mechanism or model, $\Theta_G$ is the family of claims supported by the abstraction, and $\mathcal{I}_G$ is a set of instantiation maps.

\subsection{Instantiation Maps}

A branch $c_i$ is considered an instance of $G$ if there exists an instantiation map:
$$
\psi_i:G\rightarrow c_i
$$
such that:
$$
\psi_i(O_G)\approx O_i,
$$
$$
A_i \Rightarrow \psi_i(A_G),
$$
and:
$$
\psi_i(\Theta_G)\approx H_i.
$$

When these conditions hold, we write:
$$
c_i\preceq G.
$$

This relation means that branch $c_i$ is a coherent special case, application, corollary, experiment, or boundary case of the proposed abstraction $G$.

\subsection{Abstraction Score}

The controller scores a proposed abstraction by:
$$
A(G,\mathcal{B})
=
0.25C_{\text{coverage}}
+
0.25C_{\text{coherence}}
+
0.20C_{\text{compression}}
+
0.15C_{\text{generativity}}
+
0.15C_{\text{verifiability}}.
$$

Coverage measures how many surviving branches are explained by the abstraction:
$$
C_{\text{coverage}}
=
\frac{|\{c_i\in\mathcal{B}:c_i\preceq G\}|}{|\mathcal{B}|}.
$$

Coherence measures whether the objects, assumptions, methods, and claims of the covered branches are mutually compatible. Compression measures whether $G$ simplifies the final paper rather than adding vague complexity. Generativity measures whether $G$ suggests natural new special cases, corollaries, predictions, or empirical tests. Verifiability measures whether $G$ itself, or its instances, can be supported by Lean proofs, empirical contracts, literature evidence, or explicit limitations.

A proposed abstraction is admissible only if:
$$
A(G,\mathcal{B})\geq \tau_A.
$$

The default MVP threshold is:
$$
\tau_A=0.75.
$$

\subsection{Abstraction Attack}

To prevent fake unification, every proposed abstraction is subjected to an adversarial abstraction attack. The attack asks whether the proposed general model is genuinely explanatory or merely decorative.

The abstraction attack evaluates:
\begin{itemize}[leftmargin=*]
    \item whether the covered branches are true instances of $G$;
    \item whether their assumptions are compatible;
    \item whether $G$ compresses the paper rather than adding vague terminology;
    \item whether at least one verified result depends on the abstraction;
    \item whether the abstraction improves the final paper structure.
\end{itemize}

The abstraction red-team score is:
$$
RT_{\text{abstract}}(G)\in[0,1].
$$

The abstraction survives adversarial review only if:
$$
RT_{\text{abstract}}(G)\geq \tau_{\text{abstract}}.
$$

The default MVP threshold is:
$$
\tau_{\text{abstract}}=0.75.
$$

If the abstraction fails this attack, the system must not write the final paper around $G$. It must either repair the abstraction, downgrade it to a conjectural framing, or fall back to a single verified branch.

\subsection{Final Nucleus Selection}

The final paper nucleus is selected by:
$$
\operatorname{FinalNucleus}
=
\begin{cases}
G^\star, & \text{if a valid abstraction covers multiple surviving branches},\\
c^\star, & \text{otherwise}.
\end{cases}
$$

Here:
$$
G^\star
=
\operatorname*{arg\,max}_{G\in\mathcal{G}_{\text{cand}}}
A(G,\mathcal{B}),
$$
subject to:
$$
A(G^\star,\mathcal{B})\geq\tau_A
\quad\text{and}\quad
RT_{\text{abstract}}(G^\star)\geq\tau_{\text{abstract}}.
$$

If no valid abstraction survives, the fallback branch is:
$$
c^\star
=
\operatorname*{arg\,max}_{c\in\mathcal{B}}
S_{\text{final}}(c).
$$

A default final branch score is:
$$
S_{\text{final}}(c)
=
0.30V_{\text{final}}(c)
+
0.20N(c)
+
0.20I(c)
+
0.15F(c)
+
0.10RT(c)
+
0.05R_{\text{reuse}}(c).
$$

Thus, convergence is not forced as ``many branches to one arbitrary winner.'' The system first attempts:
$$
\text{many branches}
\rightarrow
\text{shared abstraction}
\rightarrow
\text{instances}
\rightarrow
\text{one coherent paper}.
$$

Only if this fails does it use:
$$
\text{many branches}
\rightarrow
\text{best verified branch}
\rightarrow
\text{one focused paper}.
$$

\subsection{Paper Structure Induced by the Final Nucleus}

If the final nucleus is a branch $c^\star$, the manuscript is structured around a specific result:
$$
\text{Problem}
\rightarrow
\text{Method}
\rightarrow
\text{Result}
\rightarrow
\text{Verification}.
$$

If the final nucleus is an abstraction $G^\star$, the manuscript is structured around a general model:
$$
\text{General Model}
\rightarrow
\text{Instance 1}
\rightarrow
\text{Instance 2}
\rightarrow
\text{Instance 3}
\rightarrow
\text{Unified Result}
\rightarrow
\text{Limits}.
$$

In this case, surviving branches are not discarded. They become examples, special cases, corollaries, experiments, ablations, or negative boundary cases of the general model.

\section{Staged Gate Logic}

\subsection{Stage A Gate}
A candidate passes Stage A if:
$$
N(c)\geq0.35, \qquad F(c)\geq0.60, \qquad V(c)\geq0.50.
$$

\subsection{Stage B Gate}
A candidate passes Stage B if:
$$
N(c)\geq0.45, \qquad F(c)\geq0.75, \qquad V(c)\geq0.70, \qquad \sum_i f_i\leq1.
$$

\subsection{MVP Data Requirement Gate}
During the initial paper-factory implementation, a candidate may enter empirical Stage C only if:
$$
\operatorname{DataGate}_{\text{MVP}}(c)=1.
$$
Equivalently:
$$
\delta_{\text{data}}(c)\in\{\textsc{NoData},\textsc{SyntheticOnly}\}.
$$
If this gate fails, the branch is deferred rather than forced through validation:
$$
c.\text{status}\leftarrow\textsc{DeferredRealDataCandidate}.
$$

\subsection{Retrieval Adequacy Gate}
A candidate may enter Stage C only if:
$$
\rho_{\text{adequacy}}(c)\geq\tau_{\text{adequacy}}
$$
and:
$$
\rho_{\text{novelty-risk}}(c)\leq\tau_{\text{novelty-risk}}.
$$
If this gate fails, the branch receives the status:
$$
\textsc{InsufficientRetrievalAdequacy}.
$$

\subsection{Red Team Gate}
A candidate passes the Red Team Validation Layer if:
$$
RT(c)\geq0.75
$$
and its uncertainty bound holds:
$$
\hat{S}(c)-u_S(c)\geq\tau_S.
$$

\subsection{Stagnation Gate}
A branch may not remain in an unproductive repair loop. If:
$$
\operatorname{Stag}(c,t)\geq N_{\text{stag}},
$$
the orchestrator must autonomously simplify, downgrade, convert to a negative result, or stop the branch, unless $\operatorname{HumanRequired}(c,t)=1$.

\subsection{Stage C Gate}
A candidate passes Stage C if:
$$
N(c)\geq0.45, \qquad F(c)\geq0.80, \qquad V(c)\geq0.85, \qquad \sum_i f_i=0,
$$
and:
$$
\begin{cases}
V_{\text{math\_final}}(c)=1.00, & \text{if } c \text{ is a mathematical branch},\\
V_{\text{exp}}(c)=1.00 \text{ via } \textsc{StopSuccess}, & \text{if } c \text{ is an experimental branch and } \operatorname{DataGate}_{\text{MVP}}(c)=1.
\end{cases}
$$

\subsection{Abstraction Gate}
A proposed final abstraction $G$ is admissible only if:
$$
A(G,\mathcal{B})\geq\tau_A
\quad\text{and}\quad
RT_{\text{abstract}}(G)\geq\tau_{\text{abstract}}.
$$
Otherwise, the final synthesis falls back to the best verified branch $c^\star$.

\section{Human Override and Tail Escalation Policy}

A human operator may preserve a branch that fails an automatic gate only by assigning it the label $\textsc{HumanPreserved}$. It may continue, but no final claim may be made unless the corresponding final verification gate is passed.

Human override is a preservation mechanism, not a verification mechanism:
$$
\texttt{human\_override}\not\Rightarrow\textsc{LeanVerified},
$$
$$
\texttt{human\_override}\not\Rightarrow\textsc{ExperimentVerified}.
$$

The system should not request human input for ordinary controller decisions. It requests human input only when $\operatorname{HumanRequired}(c,t)=1$ under the autonomy contract. Human-facing requests must be minimal, specific, and tied to a ledgered blocking reason.

\section{Minimal Human Control Commands}

The system exposes a small set of human control commands for tail cases and explicit overrides:
$$
\begin{aligned}
\mathcal{U}_{\text{cmd}} = \{ & \texttt{stop}, \texttt{resume}, \texttt{approve}, \texttt{reject}, \texttt{approve\_split}, \\
& \texttt{reject\_split}, \texttt{revise\_task}, \texttt{request\_critic\_report}, \texttt{human\_override} \}.
\end{aligned}
$$

The command semantics are:
\begin{center}
\begin{tabular}{lll}
\toprule
Command & Scope & Effect \\
\midrule
\texttt{stop} & run or workstream & pause execution safely \\
\texttt{resume} & run or workstream & continue from latest ledger state \\
\texttt{approve} & tail proposal & allow proposed high-risk or irreversible action to proceed \\
\texttt{reject} & tail proposal & reject proposed action and request autonomous alternative \\
\texttt{approve\_split} & high-risk proof split & enqueue approved child lemmas \\
\texttt{reject\_split} & high-risk proof split & block split and request another strategy \\
\texttt{revise\_task} & workstream & modify local goal or constraints \\
\texttt{request\_critic\_report} & branch & ask for detailed objections and risks \\
\texttt{human\_override} & gate decision & preserve branch with explicit human label \\
\bottomrule
\end{tabular}
\end{center}

Every command produces a ledger commit:
$$
u_t\in\mathcal{U}_{\text{cmd}} \Rightarrow \exists n_i\in\mathcal{N}_{\text{ledger}}
$$
such that $n_i$ records the command, the previous state, the resulting state, and the user-visible reason.

\section{Research Ledger and Version Control}

The system maintains a Git-like immutable hash-linked research ledger that records every candidate transformation, controller action, verification attempt, reviewer decision, and manuscript revision. The ledger is a directed acyclic graph:
$$
\mathcal{G}_{\text{ledger}} = (\mathcal{N}_{\text{ledger}},\mathcal{E}_{\text{ledger}}).
$$

Each commit is:
$$
n_i = (h_i,h_{\text{parent}},c_i,p_i,s_i,v_i,r_i,a_i,\eta_i,t_i),
$$
where $h_i = \Hash(h_{\text{parent}},c_i,p_i,s_i,v_i,r_i,\eta_i)$.

\subsection{Commit, Diff, and Merge Rules}
Every state-changing operation (including Adaptive Controller actions) produces a commit. No candidate may be used in the final paper unless its full ancestry exists in the ledger. 

\subsection{Reproducibility, Metrology, and Redaction}
A final paper is reproducible only if the ledger contains all prompts, models, configurations, score matrices, and Lean files. Before committing, the system applies a redaction pass for credentials and private keys. 

\section{Minimal Artifact Pipeline}

The system uses a minimal artifact pipeline to prevent context bloat and to keep verification evidence separate from presentation artifacts. During search, scoring, review, proof attempts, and experiment execution, agents do not exchange full LaTeX documents. Instead, they exchange compact structured artifacts.

The internal artifact rule is:
$$
\text{internal state} = \text{JSON} + \text{Markdown} + \text{Lean files} + \text{experiment configs} + \text{small deltas}.
$$

LaTeX is reserved for final manuscript synthesis and selected user-facing reports.

\subsection{Artifact Store}

All artifacts are stored outside the orchestration runtime. The runtime may hold references to artifacts, but the artifacts themselves live in the local artifact store:
$$
\mathcal{A}_{\text{store}} = \texttt{runs/}\langle\texttt{run\_id}\rangle \cup \texttt{documents/lemmas/}.
$$

A minimal run directory is:
$$
\begin{aligned}
\texttt{runs/}\langle\texttt{run\_id}\rangle/ &= \{ \texttt{candidates/}, \texttt{scores/}, \texttt{reports/}, \texttt{literature/}, \\
&\qquad \texttt{lean/}, \texttt{experiments/}, \texttt{logs/}, \texttt{latex/} \}.
\end{aligned}
$$

Each artifact is represented by:
$$
a_i = ( \text{id}_i, \text{type}_i, \text{path}_i, \text{hash}_i, \text{created\_by}_i, \text{commit\_hash}_i, \text{metadata}_i ).
$$

The artifact hash is:
$$
\text{hash}_i=\Hash(\text{contents}(a_i)).
$$

No artifact may be used as evidence unless its hash and producing ledger commit are recorded.

\subsection{Artifact Types}

The active artifact set for a candidate branch is:
$$
A(c) = A_{\text{json}} \cup A_{\text{md}} \cup A_{\text{lean}} \cup A_{\text{exp}} \cup A_{\text{lit}} \cup A_{\text{log}} \cup A_{\text{latex}}.
$$

\begin{center}
\begin{tabular}{lll}
\toprule
Artifact Set & Format & Purpose \\
\midrule
$A_{\text{json}}$ & JSON & machine-readable states, scores, questioner reports, and controller reports \\
$A_{\text{md}}$ & Markdown & human-readable summaries and lemma pages \\
$A_{\text{lean}}$ & Lean files & formal theorem statements and proofs \\
$A_{\text{exp}}$ & configs/logs/results & experiment definitions, seeds, runs, and outputs \\
$A_{\text{lit}}$ & JSON/Markdown & retrieved papers, citation records, and coverage reports \\
$A_{\text{log}}$ & text/JSONL & controller actions, model traces, and error logs \\
$A_{\text{latex}}$ & LaTeX/PDF & final manuscript and polished reports \\
\bottomrule
\end{tabular}
\end{center}

LaTeX artifacts may not be treated as verification evidence unless linked to their underlying Lean, experiment, literature, or ledger artifacts.

\subsection{Delta-Based Communication}

Agents communicate changes rather than full histories. A state update is represented as a delta:
$$
\delta_t = ( \Delta c_t, \Delta s_t, \Delta V_t, \Delta R_t, \Delta A_t, h_t ),
$$
where $\Delta c_t$ is the candidate change, $\Delta s_t$ is the score change, $\Delta V_t$ is the verification change, $\Delta R_t$ is the review change, $\Delta A_t$ is the artifact change, and $h_t$ is the corresponding ledger commit hash.

The standard agent message is:
$$
\operatorname{Message}_t = ( \text{local context}, \delta_t, \text{requested action} ).
$$

This prevents agents from repeatedly passing full manuscripts, full branch histories, or large logs through model context.


\subsection{Runtime Context Compression}

Agents do not pass full histories through model context. The ledger keeps the full history, while runtime messages receive compressed summaries:
$$
\operatorname{RuntimeSummary}(c)=\operatorname{Compress}(\operatorname{LedgerHistory}(c)).
$$
A runtime summary may include repair counts, unresolved blockers, latest scores, current verification labels, active artifacts, and the last nontrivial improvement. It is a convenience artifact only:
$$
\operatorname{RuntimeSummary}(c)\not\Rightarrow\operatorname{Provenance}(c).
$$
If a runtime summary conflicts with the ledger, the ledger dominates.


\subsection{Evidence Boundary}

The system distinguishes presentation artifacts from evidence artifacts:
$$
A_{\text{evidence}} = A_{\text{json}} \cup A_{\text{lean}} \cup A_{\text{exp}} \cup A_{\text{lit}} \cup A_{\text{log}}.
$$

$$
A_{\text{presentation}} = A_{\text{md}} \cup A_{\text{latex}}.
$$

A final claim may cite a presentation artifact only if that presentation artifact points to evidence artifacts:
$$
\operatorname{ClaimAdmissible}(\kappa)=1 \Rightarrow \exists a_i\in A_{\text{evidence}} \text{ such that } \operatorname{Supports}(a_i,\kappa)=1.
$$

In particular:
$$
A_{\text{latex}} \not\Rightarrow \textsc{LeanVerified}
$$
and:
$$
A_{\text{latex}} \not\Rightarrow \textsc{ExperimentVerified}.
$$

\subsection{Runtime Boundary}

The artifact pipeline is independent of the orchestration runtime. A runtime such as LangGraph may route execution, pause for human approval, and resume workstreams, but it is not the source of truth.

The separation is:
$$
\text{Runtime}=\text{orchestration layer},
$$
$$
\text{Ledger}=\text{provenance layer},
$$
$$
\text{Artifacts}=\text{evidence layer},
$$
$$
\text{Schemas}=\text{contract layer}.
$$

Thus, a branch state is reproducible only if its ledger commits and artifact hashes are available:
$$
\operatorname{Reproducible}(c)=1 \Longleftrightarrow \operatorname{LedgerComplete}(c)=1 \quad\wedge\quad \operatorname{ArtifactsAvailable}(c)=1 \quad\wedge\quad \operatorname{HashesValid}(c)=1.
$$

No verified claim may depend on runtime memory alone.

\subsection{Recommended Tool Stack}

The artifact pipeline should be implemented with a conservative local-first tool stack. The first version should prioritize reproducibility, inspectability, and simple debugging over distributed execution.

\begin{center}
\begin{tabular}{p{0.18\linewidth} p{0.24\linewidth} p{0.48\linewidth}}
\toprule
Layer & Recommended Tool & Role in the System \\
\midrule
Core language & Python & Main implementation language for controllers, adapters, scoring, and orchestration. \\
Schema validation & Pydantic & Defines strict contracts for candidates, scores, ledger commits, controller actions, theorem records, and artifact references. \\
MVP runtime & Plain Python state machine & Executes the first deterministic version before adding graph orchestration. \\
Graph runtime & LangGraph & Added after the MVP loop is stable, for pause/resume, human approval, branching, retries, and controller routing. \\
Ledger database & SQLite & Local append-only store for commits, branch ancestry, command history, and run metadata. \\
Ledger database, later & PostgreSQL & Optional replacement when runs become large, concurrent, or multi-user. \\
Artifact store & Local filesystem & Stores JSON, markdown, Lean files, logs, experiment outputs, literature records, and LaTeX artifacts. \\
CLI & Typer or equivalent & Provides local commands such as \texttt{run}, \texttt{status}, \texttt{report}, \texttt{approve}, and \texttt{approve\_split}. \\
Testing & pytest & Validates schemas, ledger invariants, scoring, controllers, adapters, and full pipeline runs. \\
Linting and formatting & Ruff & Maintains code quality and consistent formatting. \\
Type checking & mypy or pyright & Optional static checking for larger modules and public interfaces. \\
Model adapter & Provider-agnostic wrapper & Encapsulates calls to OpenAI, Anthropic, local models, or mock models. \\
Literature adapter & Semantic Scholar, OpenAlex, arXiv, or local corpus & Retrieves prior work, novelty evidence, citation metadata, and coverage reports. \\
Vector search, later & FAISS or pgvector & Supports semantic retrieval over artifacts, papers, lemma pages, and candidate histories. \\
Formal verification & Lean 4 subprocess & Checks theorem statements and proofs in a sandboxed Lean project. \\
Experiment runtime, later & Docker sandbox & Runs empirical code with pinned dependencies, seeds, metrics, and output hashes. \\
API, later & FastAPI & Exposes run control and artifact inspection after the CLI is stable. \\
Frontend, later & Minimal dashboard & Displays run status, workstreams, approval requests, reports, and artifact links. \\
\bottomrule
\end{tabular}
\end{center}

The initial implementation stack is therefore:
$$
\text{Python} + \text{Pydantic} + \text{SQLite} + \text{local filesystem} + \text{Typer} + \text{pytest} + \text{Ruff}.
$$

LangGraph, Lean, vector search, Docker experiments, FastAPI, and a dashboard are added only after the local deterministic pipeline passes its version gates.

\subsection{Tool Boundary Rules}

The implementation must preserve the following boundaries:
$$
\text{Pydantic}=\text{contract validation},
$$
$$
\text{SQLite}=\text{ledger metadata},
$$
$$
\text{filesystem}=\text{artifact storage},
$$
$$
\text{LangGraph}=\text{orchestration},
$$
$$
\text{Lean}=\text{formal verification},
$$
$$
\text{Docker}=\text{empirical sandboxing}.
$$

No tool is allowed to replace the ledger as the provenance layer:
$$
\text{Runtime state} \not\Rightarrow \text{provenance}.
$$

No tool is allowed to replace verification artifacts:
$$
\text{model output} \not\Rightarrow \textsc{LeanVerified},
$$
$$
\text{model output} \not\Rightarrow \textsc{ExperimentVerified}.
$$

The graph runtime may decide where execution goes next, but it may not silently mutate candidate state. Every runtime transition that changes a candidate, score, artifact, verification label, or branch status must produce a ledger commit.

\subsection{Tool Introduction Order}

Tools are introduced in the following order:
$$
\begin{aligned}
&1.\ \text{Python, Pydantic, pytest, Ruff},\\
&2.\ \text{SQLite ledger and local filesystem artifact store},\\
&3.\ \text{Typer CLI},\\
&4.\ \text{fake model, fake literature, and fake embedding adapters},\\
&5.\ \text{real model adapter with schema validation},\\
&6.\ \text{real literature adapter},\\
&7.\ \text{LangGraph orchestration},\\
&8.\ \text{Lean 4 subprocess and theorem registry},\\
&9.\ \text{vector search over artifacts and lemma pages},\\
&10.\ \text{Docker experiment sandbox},\\
&11.\ \text{FastAPI and lightweight dashboard}.
\end{aligned}
$$

The system may introduce tool $T_{i+1}$ only if tool $T_i$ is covered by tests and does not bypass the ledger or artifact invariants.

\subsection{Minimum Artifact Invariants}

The implementation must enforce the following invariants:
$$
\begin{aligned}
&\textbf{I1:} \quad \text{Every state-changing action creates a ledger commit.}\\
&\textbf{I2:} \quad \text{Every artifact has a content hash.}\\
&\textbf{I3:} \quad \text{Every artifact is linked to a producing commit.}\\
&\textbf{I4:} \quad \text{Every final claim links to evidence artifacts.}\\
&\textbf{I5:} \quad \text{No presentation artifact counts as verification evidence.}\\
&\textbf{I6:} \quad \text{No runtime state alone counts as provenance.}
\end{aligned}
$$

A build fails if any invariant is violated.



\section{Global Search Execution Pipeline}

To prevent monolithic logic blocks and page overflow, the execution pipeline is decomposed into page-safe modular sub-routines. The top-level orchestrator manages transitions between routines while the ledger and artifact store remain the sources of truth. The pipeline is autonomous by default: human intervention occurs only when $\operatorname{HumanRequired}(c,t)=1$.

\begin{algorithm}[H]
\caption{Factory-Mode Master Orchestration Pipeline}
\begin{algorithmic}[1]
\Require User constraint set $C$, method candidate source $\mathcal{M}_{\text{cand}}$, budget vector $\mathbf{B}$
\Ensure Verified manuscript, negative-result report, conjecture report, or failed-search report
\State Initialize run id $\rho$, artifact store $\mathcal{A}_{\text{store}}$, and ledger $\mathcal{G}_{\text{ledger}}$
\State Commit root constraint set $C$
\State $\Cands_A\gets$ \Call{Stage0\_and\_StageA}{$C,\mathcal{M}_{\text{cand}}$}
\State $\Cands_B\gets$ \Call{StageB\_StructuralValidation}{$\Cands_A,\mathbf{B}$}
\State $\Cands_C\gets$ \Call{RedTeam\_Filter}{$\Cands_B$}
\State $\mathcal{B}\gets$ \Call{StageC\_DeepVerification}{$\Cands_C,\mathbf{B}$}
\If{$\mathcal{B}\neq\emptyset$}
    \State $N_{\text{final}}\gets$ \Call{AbstractSynthesis\_and\_FinalSelection}{$\mathcal{B}$}
    \State \Return Compile labeled manuscript and evidence bundle from $N_{\text{final}}$
\Else
    \State \Return Emit failed-search report with ledgered reasons and artifacts
\EndIf
\end{algorithmic}
\end{algorithm}

\begin{algorithm}[H]
\footnotesize
\caption{Opportunity Discovery and Stage A Expansion}
\begin{algorithmic}[1]
\Function{Stage0\_and\_StageA}{$C,\mathcal{M}_{\text{cand}}$}
\If{$C$ lacks a specified method $M$ but defines a domain $D$}
    \State Extract primitives $P(D)$; write artifact and commit it
    \For{each method $M_i\in\mathcal{M}_{\text{cand}}$}
        \State Compute $O_{\text{final}}(D,M_i)$; write score artifact and commit it
    \EndFor
    \State Construct $\mathcal{C}_{0}\gets\{C\cup(M=M_i):O_{\text{final}}(D,M_i)\geq\tau_{\text{opportunity}}\}$
\Else
    \State Set $\mathcal{C}_{0}\gets\{C\}$
\EndIf
\State Generate candidate branches satisfying seeded constraints
\State Validate schemas, write candidate artifacts, and commit candidate creation
\For{each candidate $c$}
    \State Run Question Refiner and Retrieval Adequacy Controller; commit actions
    \State Compute scores and write score artifact
    \State Select $Q_t$ and run Strategic Questioner; route autonomous action if needed
    \State Update $\operatorname{Stag}(c,t)$ and resolve stagnant branches autonomously
\EndFor
\State Apply duplicate filtering, Stage A gate, and top-$K$ survivor selection
\State Commit Stage A survivor set and return it
\EndFunction
\end{algorithmic}
\end{algorithm}

\begin{algorithm}[H]
\footnotesize
\caption{Stage B Structural Validation}
\begin{algorithmic}[1]
\Function{StageB\_StructuralValidation}{$\Cands_A,\mathbf{B}$}
\State Expand survivors into localized child variants; validate and commit children
\State Execute Reviewer Disagreement Resolver; write review artifacts and commit decisions
\For{each reviewed candidate $c$}
    \State Select $Q_t$ and run Strategic Questioner
    \If{meta-action is autonomous and blocking}
        \State Route to the selected controller before Stage B gating
    \ElsIf{$\operatorname{HumanRequired}(c,t)=1$}
        \State Emit minimal human steering packet and commit $\textsc{HumanTailEscalation}$
    \EndIf
    \State Update $\operatorname{Stag}(c,t)$ and force resolution if needed
\EndFor
\State Apply Stage B gate, rank by $S_{\text{cost}}(c)$, and allocate budget
\State Commit Stage B survivor set and return it
\EndFunction
\end{algorithmic}
\end{algorithm}

\begin{algorithm}[H]
\footnotesize
\caption{Red Team Filtering}
\begin{algorithmic}[1]
\Function{RedTeam\_Filter}{$\Cands_B$}
\State Initialize $\Cands_C\gets\emptyset$
\For{each Stage B survivor $c$}
    \State Execute novelty, bridge, baseline, and triviality attacks
    \State Check MVP data gate $\operatorname{DataGate}_{\text{MVP}}(c)$ and retrieval adequacy certificate $\rho_{\text{adequacy}}(c)$
    \If{$\operatorname{DataGate}_{\text{MVP}}(c)=0$}
        \State Mark $c$ as $\textsc{DeferredRealDataCandidate}$ and commit
    \ElsIf{$\rho_{\text{adequacy}}(c)<\tau_{\text{adequacy}}$}
        \State Mark $c$ as $\textsc{InsufficientRetrievalAdequacy}$ and commit
    \ElsIf{$RT_{\text{bridge}}(c)=0$}
        \State Execute one autonomous Bridge Controller repair attempt; commit repair
    \EndIf
    \State Compute $RT(c)$ and conservative bound $\hat{S}(c)-u_S(c)$
    \State Select $Q_t$ and run Strategic Questioner; route autonomous repair if selected
    \If{$RT(c)<\tau_{\text{redteam}}$}
        \State Mark $c$ as $\textsc{RejectedRedTeam}$ and commit
    \ElsIf{$\hat{S}(c)-u_S(c)<\tau_S$}
        \State Mark $c$ as $\textsc{PrunedUncertain}$ and commit
    \Else
        \State Add $c$ to $\Cands_C$
    \EndIf
\EndFor
\State Select Stage C candidates by $S_{\text{cost}}(c)$ and remaining budget
\State Commit Stage C candidate set and return it
\EndFunction
\end{algorithmic}
\end{algorithm}

\begin{algorithm}[H]
\footnotesize
\caption{Stage C Deep Verification}
\begin{algorithmic}[1]
\Function{StageC\_DeepVerification}{$\Cands_C,\mathbf{B}$}
\State Initialize verified-or-labeled branch set $\mathcal{B}\gets\emptyset$
\For{each Stage C candidate $c$}
    \If{$c$ is mathematical}
        \State Run Adaptive Proof Controller and Lean verification
        \State Apply bounded auto proof splitting when allowed
        \State Escalate only if $\operatorname{HumanRequired}(c,t)=1$
        \State Assign $\textsc{LeanVerified}$, $\textsc{Conjecture}$, or $\textsc{LeanFailure}$
    \ElsIf{$c$ is empirical}
        \State Apply MVP data gate and run Adaptive Experiment Controller in sandbox
        \State Assign $\textsc{SyntheticExperimentVerified}$, $\textsc{NegativeResult}$, $\textsc{DeferredRealDataCandidate}$, or $\textsc{ExperimentFailure}$
    \Else
        \State Mark unsupported branch type as $\textsc{Limitation}$ or reject it
    \EndIf
    \State Select $Q_t$ and run Strategic Questioner; autonomously downgrade or stop if needed
    \State Update $\operatorname{Stag}(c,t)$ and force resolution if stalled
    \If{$c$ has an admissible final label}
        \State Add $c$ to $\mathcal{B}$ and commit final branch label
    \EndIf
\EndFor
\State \Return $\mathcal{B}$
\EndFunction
\end{algorithmic}
\end{algorithm}

\begin{algorithm}[H]
\footnotesize
\caption{Abstract Synthesis and Final Output Selection}
\begin{algorithmic}[1]
\Function{AbstractSynthesis\_and\_FinalSelection}{$\mathcal{B}$}
\State Generate candidate abstractions $\mathcal{G}_{\text{cand}}=\{G_1,\dots,G_m\}$
\For{each abstraction $G_j$}
    \State Construct instantiation maps $\psi_i:G_j\rightarrow c_i$
    \State Compute $A(G_j,\mathcal{B})$ and run Abstraction Attack
    \State Commit abstraction report and risk score
\EndFor
\State Let $G^\star=\operatorname*{arg\,max}_{G_j}A(G_j,\mathcal{B})$
\If{$A(G^\star,\mathcal{B})\geq\tau_A$ and $RT_{\text{abstract}}(G^\star)\geq\tau_{\text{abstract}}$}
    \State Set $N_{\text{final}}\gets G^\star$ and commit $\textsc{AbstractNucleus}$
\Else
    \State Set $c^\star=\operatorname*{arg\,max}_{c\in\mathcal{B}}S_{\text{final}}(c)$
    \State Set $N_{\text{final}}\gets c^\star$ and commit $\textsc{BranchNucleus}$
\EndIf
\If{two final nuclei are tied within $\epsilon_{\text{tie}}$ and $\operatorname{HumanRequired}=1$}
    \State Emit minimal human preference packet
\EndIf
\State Link every final claim to supporting evidence artifacts
\State Write markdown report, LaTeX source, and PDF artifact; commit final synthesis
\State \Return $N_{\text{final}}$
\EndFunction
\end{algorithmic}
\end{algorithm}

\section{Synthesis Specifications and Final Output Schema}

Claims in the manuscript must be labeled as:
$$
\textsc{LeanVerified}, \quad
\textsc{SyntheticExperimentVerified}, \quad
\textsc{RealDataExperimentVerified}, \quad
\textsc{Conjecture}, \quad
\textsc{NegativeResult}, \quad
\textsc{Limitation}.
$$
For compactness, \textsc{ExperimentVerified} denotes either \textsc{SyntheticExperimentVerified} or \textsc{RealDataExperimentVerified}, but the two labels are not interchangeable.

No theorem declaration may be generated unless $V_{\text{math\_final}}(c)=1.00$. Unresolved mathematical structures at Lean levels $1$ through $4$ must be labeled $\textsc{Conjecture}$. No final empirical claim may be labeled $\textsc{SyntheticExperimentVerified}$ or $\textsc{RealDataExperimentVerified}$ unless the experiment controller reaches $\textsc{StopSuccess}$ under a predeclared experimental contract. Under MVP factory constraints, empirical success based on synthetic data must be labeled $\textsc{SyntheticExperimentVerified}$ and may not be presented as real-world validation.

The final manuscript is organized around a final nucleus:
$$
N_{\text{final}}\in\{G^\star,c^\star\}.
$$

If $N_{\text{final}}=G^\star$, the manuscript is written as a general model paper:
$$
\text{General Model}
\rightarrow
\text{Instances}
\rightarrow
\text{Verified Results}
\rightarrow
\text{Boundary Cases}
\rightarrow
\text{Limitations}.
$$

If $N_{\text{final}}=c^\star$, the manuscript is written as a focused branch paper:
$$
\text{Problem}
\rightarrow
\text{Method}
\rightarrow
\text{Result}
\rightarrow
\text{Verification}
\rightarrow
\text{Limitations}.
$$

The final output is:
$$
O=
(P_{\text{final}},
N_{\text{final}},
\mathcal{G}_{\text{ledger}},
\mathcal{B}_{\text{secondary}},
\mathcal{B}_{\text{failed}},
\mathcal{A}_{\text{evidence}},
\mathcal{R}).
$$

Here $P_{\text{final}}$ is the labeled manuscript, $N_{\text{final}}$ is the final nucleus, $\mathcal{B}_{\text{secondary}}$ contains surviving branches not used as the main nucleus, $\mathcal{B}_{\text{failed}}$ contains rejected or failed branches, $\mathcal{A}_{\text{evidence}}$ contains proof, experiment, literature, score, and red-team artifacts, and $\mathcal{R}$ contains review and synthesis reports.

If all branches fail final gates, the system returns a failed-search report with reason codes:
$$
\begin{aligned}
\text{Failure Reason} \in \{ &\textsc{LowNovelty}, \textsc{LowFeasibility}, \textsc{LowVerifiability}, \\
&\textsc{LeanFailure}, \textsc{ExperimentFailure}, \textsc{ReviewerFatalFlaw}, \\
&\textsc{BudgetExhausted}, \textsc{DuplicateBranch}, \textsc{FalseBridge}, \\
&\textsc{RejectedRedTeam}, \textsc{PrunedUncertain}, \textsc{InvalidAbstraction}, \\
&\textsc{InsufficientRetrievalAdequacy}, \textsc{StagnationStop}, \\
&\textsc{AutoSplitRejected}, \textsc{HumanTailEscalation}, \\
&\textsc{AbstractionRejected}, \textsc{CalibrationFlagged}, \\
&\textsc{RequiresRealData}, \textsc{DeferredRealDataCandidate}, \\
&\textsc{AutonomyBlocked}, \textsc{NonReproducible} \}.
\end{aligned}
$$

The synthesis stage enforces output normalization rather than unconditional paper generation:
$$
\operatorname{Output}(\Tree)
=
\begin{cases}
\textsc{VerifiedManuscript}, & \text{if a verified or valid labeled nucleus exists},\\
\textsc{NegativeResultReport}, & \text{if only informative negative results survive},\\
\textsc{FailedSearchReport}, & \text{if no admissible final nucleus survives}.
\end{cases}
$$

\appendix
\section{Research Density Landscape and Controller-Guided Descent}
This appendix gives an interpretation of the search procedure as a derivative-free descent process over a latent research landscape. The purpose is not to claim that scientific discovery is a smooth optimization problem, but to make explicit the objective implicitly used by the staged search, controller refinements, and pruning rules.

\subsection{Admissible Paper Space}
Let $C$ be the user constraint set. A minimal admissible paper description is written as:
$$ x = (D,P,M,Q,H,T,E,B), $$
where $D$ is the domain, $P$ the extracted primitives, $M$ the method, $Q$ the research question, $H$ the hypothesis, $T$ the theoretical structure, $E$ the experiment contract, and $B$ the baseline.

The admissible research space is defined as:
$$ \mathcal{X}_C = \{x : x \models C\}. $$

Each $x\in\mathcal{X}_C$ is mapped into a latent research space by an ensemble embedding:
$$ \phi:\mathcal{X}_C\rightarrow \mathbb{R}^d. $$

Distances in this space are estimated using the framework's ensemble branch distance:
$$ d_{\mathcal{X}}(x_i,x_j)=D_{\mathrm{track}}(x_i,x_j). $$

\subsection{Prior-Work Density}
Let $\mathcal{P}$ be the retrieved literature corpus. The local prior-work density around $x$ is estimated by:
$$ \rho_{\mathrm{lit}}(x) = \sum_{p\in\mathcal{P}} \exp\left( -\frac{d_{\mathcal{X}}(x,p)^2}{2\sigma^2} \right). $$

High $\rho_{\mathrm{lit}}(x)$ indicates a crowded region of prior work. Low $\rho_{\mathrm{lit}}(x)$ indicates a sparse or underexplored region. The density estimate is considered valid only when the literature controller has sufficient retrieval adequacy:
$$ \rho_{\mathrm{adequacy}}(x)\geq \tau_{\mathrm{adequacy}}. $$

\subsection{Difficulty-Weighted Energy}
Low density alone is not sufficient, since a region may be sparse because it is impossible, irrelevant, or unverifiable. The system therefore defines an energy function:
$$ \mathcal{E}(x) = \rho_{\mathrm{lit}}(x) + \lambda_{\mathrm{diff}}\operatorname{Diff}(x) + \lambda_{\mathrm{risk}}R_{\mathrm{risk}}(x) - \lambda_{\mathrm{val}}I(x)F(x)V(x), $$
where $\operatorname{Diff}(x)$ is difficulty, $R_{\mathrm{risk}}(x)$ aggregates novelty, bridge, baseline, and verification risk, and $I(x)$, $F(x)$, and $V(x)$ denote importance, feasibility, and verifiability.

The system seeks low-energy regions:
$$ x^\star \in \operatorname*{arg\,locmin}_{x\in\mathcal{X}_C} \mathcal{E}(x). $$

A good candidate is therefore not merely novel. It is sparse in prior work, feasible, verifiable, scientifically valuable, and not too difficult.

\subsection{Controller-Guided Local Search}
The system does not perform classical gradient descent. The variables in $x$ are partly semantic, partly symbolic, and partly discrete. Therefore a true gradient $\nabla \mathcal{E}(x)$ is generally unavailable.

Instead, the research tree performs derivative-free, controller-guided local search. Each adaptive controller proposes neighboring candidates:
$$ \mathcal{N}(x_t) = \{a(x_t):a\in\mathcal{A}_{\mathrm{ctrl}}\}, $$
where the control actions break down as:
$$ 
\begin{aligned}
\mathcal{A}_{\mathrm{ctrl}} = \{ &\textsc{NarrowScope}, \textsc{ChangeMetric}, \textsc{AddBaseline}, \textsc{SearchAdjacentFields}, \\
&\textsc{DefineObjects}, \textsc{AddDataset}, \textsc{StrengthenBaseline}, \dots \}.
\end{aligned}
$$

A local move is selected by approximately minimizing the estimated energy over the controller-generated neighborhood:
$$ x_{t+1} \approx \operatorname*{arg\,min}_{y\in\mathcal{N}(x_t)} \widehat{\mathcal{E}}(y). $$

The move is accepted if it improves the conservative objective:
$$ \widehat{\mathcal{E}}(x_{t+1})+\eta_t < \widehat{\mathcal{E}}(x_t), $$
where $\eta_t$ is an exploration tolerance used to preserve promising high-variance branches.

\subsection{Underexploited Feasible Basins}
A region $U\subset\mathcal{X}_C$ is an underexploited feasible basin if:
$$ \mathbb{E}_{x\in U}[\rho_{\mathrm{lit}}(x)]\leq \tau_\rho, $$
$$ \tau_{\mathrm{diff}, \min} \leq \mathbb{E}_{x\in U}[\operatorname{Diff}(x)] \leq \tau_{\mathrm{diff},\max}, $$
and:
$$ \mathbb{E}_{x\in U}[F(x)V(x)]\geq \tau_{FV}. $$

Stage 0 Opportunity Discovery can be interpreted as a coarse basin search. For a domain-method pair $(D,M)$, define:
$$ U_{D,M} = \{x\in\mathcal{X}_C:x_D=D,\ x_M=M\}. $$

Then the opportunity score approximates a basin score:
$$ O_{\mathrm{final}}(D,M) \approx \operatorname{BasinScore}(U_{D,M}). $$

\subsection{Interpretation}
The resulting search dynamics are best described as a \textit{controller-guided descent over a difficulty-weighted research density landscape}.

This interpretation explains why the framework maintains variance early and applies stronger validation later. Early stages sample and explore possible basins; middle stages refine and compare local neighborhoods; final stages verify only the surviving low-energy candidates. The final manuscript is not the landscape itself, but the projection of a surviving candidate into a labeled scientific artifact with explicit provenance, evidence, and verification status.
\end{document}
